\documentclass[11pt,a4paper]{article}

% --- PAQUETES FUNDAMENTALES ---
\usepackage[utf8]{inputenc}
\usepackage[spanish,es-nodecimaldot,es-noshorthands]{babel}
\usepackage{amsmath,amssymb,amsfonts,amsthm,mathtools}
\usepackage{geometry}
\geometry{top=2cm, bottom=2cm, left=2cm, right=2cm}
\usepackage{xcolor}
\usepackage{tcolorbox}
\tcbuselibrary{skins,breakable}
\usepackage{fancyhdr}
\usepackage{hyperref}
\usepackage{enumitem}
\usepackage{tikz}
\usetikzlibrary{arrows.meta,calc,decorations.markings,patterns,babel}

% --- CONFIGURACIÓN DE PÁGINA ---
\pagestyle{fancy}
\fancyhf{}
\fancyhead[L]{\small \textbf{UNMSM -- FCM}}
\fancyhead[R]{\small \textbf{Problemario: Sistemas Dinámicos}}
\fancyfoot[C]{\thepage}
\renewcommand{\headrulewidth}{0.8pt}
\renewcommand{\footrulewidth}{0.4pt}
\raggedbottom

% --- PALETA DE COLORES ---
\definecolor{unmsmblue}{RGB}{10, 45, 95}
\definecolor{deepteal}{RGB}{18, 98, 115}
\definecolor{darkpurple}{RGB}{75, 25, 105}
\definecolor{solgreen}{RGB}{242, 249, 244}
\definecolor{solborder}{RGB}{35, 120, 75}
\definecolor{probhead}{RGB}{240, 244, 250}
\definecolor{warnbg}{RGB}{254, 245, 235}
\definecolor{noteborder}{RGB}{190, 85, 20}

% --- ENTORNO DE PROBLEMA (RECUADRO AZUL) Y SOLUCIONES LIMPIAS ---
\newcommand{\tcbbreak}{}

\newtcolorbox{problema}[1][]{
    colback=probhead,
    colframe=unmsmblue,
    fonttitle=\bfseries\large\color{white},
    title={#1},
    arc=1.5mm,
    boxrule=0.9pt,
    left=3mm, right=3mm, top=2mm, bottom=2mm,
    before skip=3.5mm,
    after skip=2.5mm,
    breakable
}

\newenvironment{formulas}[1][Marco Teórico y Fórmulas Empleadas]{%
    \vspace{2mm}%
    \noindent{\large\textbf{\textit{#1:}}}\par\vspace{1.5mm}%
}{%
    \par\vspace{2mm}%
}

\newenvironment{solucion}[1][Solución Detallada]{%
    \vspace{1.5mm}%
    \noindent{\large\textbf{Solución:}}\par\vspace{1.5mm}%
}{%
    \par\vspace{3mm}%
}

\newenvironment{conclusion}[1][Interpretación Dinámica y Conclusión]{%
    \vspace{2mm}%
    \noindent{\large\textbf{Conclusión:}}\par\vspace{1.5mm}%
}{%
    \par\vspace{3mm}%
}

\hypersetup{
    colorlinks=true,
    linkcolor=unmsmblue,
    urlcolor=deepteal,
    citecolor=unmsmblue
}

\begin{document}

% --- PORTADA ---
\begin{titlepage}
\begin{center}
    \vspace*{2.0cm}
    
    {\fontsize{26}{32}\selectfont \textbf{Problemario de Ecuaciones Diferenciales}}\\[3mm]
    {\fontsize{26}{32}\selectfont \textbf{Ordinarias}}\\[8mm]
    {\LARGE \textbf{Grupo 09}}\\[2.0cm]
    
    {\normalsize Illescas Vicente Alexander George, \quad Tisnado Yarleque Christian David,}\\[2.5mm]
    {\normalsize Yanac Minaya Junior Alberto, \quad Quispe Gonzales Mark, \quad David Alejandro Tejada Ossio}\\[2.2cm]
    
    \renewcommand{\arraystretch}{1.35}
    \begin{tabular}{rl}
        \textbf{Profesor} & : JULIOS DECIDERIO CASTILLO MELCHOR \\
        \textbf{Curso} & : Redes, Arquitectura y Comunicaciones \\
        \textbf{Escuela Profesional} & : Computación Científica
    \end{tabular}\\[2.5cm]
    
    \vfill
    
    {\large \textbf{Universidad Nacional Mayor de San Marcos}}\\[2mm]
    {\normalsize Lima, Perú -- 2026}
    
    \vspace*{1.0cm}
\end{center}
\end{titlepage}

% =========================================================================
% TEMA 1: EDOs LINEALES Y NO LINEALES
% =========================================================================
\section*{\color{unmsmblue} Tema 1: EDOs Lineales y No Lineales}

\begin{problema}[Problema 1.1]
Resuelva la ecuación diferencial ordinaria separable:
\[
\frac{dy}{dx} = 3x y^2, \qquad y(0) = 1
\]
y determine el intervalo máximo de existencia de la solución.
\end{problema}

\begin{solucion}

\textbf{\large 1. Separar variables}
\[
\frac{dy}{dx} = 3x y^2
\]
\[
\frac{dy}{y^2} = 3x\, dx
\]
\[
\boxed{y^{-2}\, dy = 3x\, dx}
\]

\vspace{3.5mm}

\textbf{\large 2. Integrar ambos miembros}
\[
\int y^{-2}\, dy = \int 3x\, dx
\]
\[
-\frac{1}{y} = 3\left(\frac{x^2}{2}\right) + C
\]
\[
\boxed{-\frac{1}{y} = \frac{3}{2}x^2 + C}
\]

\vspace{3.5mm}

\textbf{\large 3. Despejar la variable $y(x)$}
\[
\frac{1}{y} = -\left(\frac{3}{2}x^2 + C\right) = -\frac{3x^2 + 2C}{2}
\]
\[
\boxed{y(x) = -\frac{2}{3x^2 + 2C}}
\]

\vspace{3.5mm}

\textbf{\large 4. Aplicar la condición inicial $y(0) = 1$}
\[
1 = -\frac{2}{3(0)^2 + 2C}
\]
\[
1 = -\frac{2}{2C} = -\frac{1}{C}
\]
\[
\boxed{C = -1}
\]
\[
y(x) = -\frac{2}{3x^2 + 2(-1)}
\]
\[
\boxed{y(x) = \frac{2}{2 - 3x^2}}
\]

\vspace{3.5mm}

\textbf{\large 5. Determinar el intervalo maximal de existencia}
\[
2 - 3x^2 = 0 \implies x^2 = \frac{2}{3} \implies x = \pm\sqrt{\frac{2}{3}}
\]
\[
x_0 = 0 \in \left(-\sqrt{\frac{2}{3}}, \, \sqrt{\frac{2}{3}}\right)
\]
\[
\boxed{I = \left(-\sqrt{\frac{2}{3}}, \, \sqrt{\frac{2}{3}}\right)}
\]

\begin{center}
\begin{tikzpicture}[scale=1.1]
    \draw[->, thick] (-2.2,0) -- (2.2,0) node[right] {$x$};
    \draw[->, thick] (0,-0.5) -- (0,4.2) node[above] {$y(x)$};
    \draw[dashed, red!80!black, thick] (-0.8165,-0.5) -- (-0.8165,4.0) node[pos=0.08, left, font=\scriptsize] {$x = -\sqrt{2/3}$};
    \draw[dashed, red!80!black, thick] (0.8165,-0.5) -- (0.8165,4.0) node[pos=0.08, right, font=\scriptsize] {$x = \sqrt{2/3}$};
    \draw[very thick, deepteal, domain=-0.75:0.75, samples=18, smooth] plot (\x, {2 / (2 - 3*\x*\x)});
    \filldraw[unmsmblue] (0,1) circle (2.2pt) node[below right, font=\footnotesize] {$(0, 1)$};
    \node[above, font=\footnotesize, deepteal] at (0, 1.3) {$y(x) = \frac{2}{2 - 3x^2}$};
\end{tikzpicture}
\end{center}

\end{solucion}

\vspace{3mm}

\begin{problema}[Problema 1.2]
Resuelva la ecuación diferencial ordinaria de Bernoulli:
\[
y' + y = y^3, \qquad y(0) = 1
\]
\end{problema}

\begin{solucion}

\textbf{\large 1. Identificar la ecuación de Bernoulli}
\[
\boxed{y' + P(x)y = Q(x)y^n}
\]
\[
y' + y = y^3
\]
\[
P(x) = 1, \qquad Q(x) = 1, \qquad n = 3
\]

\vspace{3.5mm}

\textbf{\large 2. Aplicar el cambio de variable}
\[
\boxed{v = y^{1-n}}
\]
\[
v = y^{1-3}
\]
\[
\boxed{v = y^{-2} = \frac{1}{y^2}}
\]
\[
v = y^{-2}
\]
\[
\frac{dv}{dx} = -2y^{-3} \frac{dy}{dx}
\]
\[
\boxed{v' = -2y^{-3}y'}
\]
\[
\boxed{y^{-3}y' = -\frac{1}{2}v'}
\]

\vspace{3.5mm}

\textbf{\large 3. Dividir toda la ecuación entre $y^3$}
\[
y' + y = y^3
\]
\[
\frac{y'}{y^3} + \frac{y}{y^3} = \frac{y^3}{y^3}
\]
\[
y^{-3}y' + y^{-2} = 1
\]
\[
y^{-3}y' = -\frac{1}{2}v', \qquad y^{-2} = v
\]
\[
-\frac{1}{2}v' + v = 1
\]

\vspace{3.5mm}

\textbf{\large 4. Convertirla en una ecuación lineal}
\[
-\frac{1}{2}v' + v = 1
\]
\[
\boxed{v' - 2v = -2}
\]

\vspace{3.5mm}

\textbf{\large 5. Resolver la ecuación lineal}
\[
v' - 2v = -2
\]
\[
v' + P(x)v = Q(x)
\]
\[
P(x) = -2
\]
\[
\mu(x) = e^{\int P(x)\,dx} = e^{\int -2\,dx}
\]
\[
\boxed{\mu(x) = e^{-2x}}
\]

\vspace{3.5mm}

\textbf{\large 6. Multiplicar por el factor integrante e integrar}
\[
e^{-2x}v' - 2e^{-2x}v = -2e^{-2x}
\]
\[
\frac{d}{dx}\left[e^{-2x}v\right] = -2e^{-2x}
\]
\[
\int \frac{d}{dx}\left[e^{-2x}v\right]\,dx = \int -2e^{-2x}\,dx
\]
\[
e^{-2x}v = -2\left(\frac{e^{-2x}}{-2}\right) + C = e^{-2x} + C
\]
\[
\boxed{v = 1 + C e^{2x}}
\]

\vspace{3.5mm}

\textbf{\large 7. Volver a la variable original $y$}
\[
v = \frac{1}{y^2}
\]
\[
\frac{1}{y^2} = 1 + C e^{2x}
\]
\[
y^2 = \frac{1}{1 + C e^{2x}}
\]
\[
y = \pm \frac{1}{\sqrt{1 + C e^{2x}}}
\]
\[
\boxed{y(x) = \frac{1}{\sqrt{1 + C e^{2x}}}}
\]

\vspace{3.5mm}

\textbf{\large 8. Aplicar la condición inicial $y(0) = 1$}
\[
1 = \frac{1}{\sqrt{1 + C e^{2(0)}}} = \frac{1}{\sqrt{1 + C}}
\]
\[
1 = \frac{1}{1 + C} \implies 1 + C = 1
\]
\[
\boxed{C = 0}
\]
\[
y(x) = \frac{1}{\sqrt{1 + 0}}
\]
\[
\boxed{y(x) \equiv 1, \qquad \forall x \in \mathbb{R}}
\]

\begin{center}
\begin{tikzpicture}[scale=1.1]
    \draw[->, thick] (-2.5,0) -- (2.5,0) node[right] {$x$};
    \draw[->, thick] (0,-0.5) -- (0,2.2) node[above] {$y(x)$};
    % Solución constante de equilibrio y = 1
    \draw[very thick, deepteal] (-2.3,1) -- (2.3,1) node[right, font=\footnotesize] {$y(x) \equiv 1$};
    \filldraw[unmsmblue] (0,1) circle (2.2pt) node[above left, font=\footnotesize] {$(0, 1)$};
\end{tikzpicture}
\end{center}

\end{solucion}

\vspace{3mm}

\begin{problema}[Problema 1.3]
Halle la solución general de la ecuación diferencial:
\[
x^2 y'' - 2x y' + 2y = x^3 \ln x, \quad x > 0
\]
\end{problema}


\begin{solucion}

\textbf{\large 1. Identificar la ecuación de Cauchy-Euler}
\[
\boxed{a x^2 y'' + b x y' + c y = f(x)}
\]
\[
x^2 y'' - 2x y' + 2y = x^3 \ln x
\]
\[
a = 1, \qquad b = -2, \qquad c = 2
\]

\vspace{3.5mm}

\textbf{\large 2. Resolver la ecuación homogénea asociada}
\[
x^2 y'' - 2x y' + 2y = 0
\]
\[
y = x^m \implies y' = m x^{m-1} \implies y'' = m(m-1)x^{m-2}
\]
\[
x^2 [m(m-1)x^{m-2}] - 2x [m x^{m-1}] + 2 [x^m] = 0
\]
\[
m(m-1) - 2m + 2 = 0
\]
\[
m^2 - 3m + 2 = 0
\]
\[
(m - 1)(m - 2) = 0
\]
\[
m_1 = 1, \qquad m_2 = 2
\]
\[
y_1(x) = x, \qquad y_2(x) = x^2
\]
\[
\boxed{y_h(x) = c_1 x + c_2 x^2}
\]

\vspace{3.5mm}

\textbf{\large 3. Normalizar la ecuación diferencial}
\[
y'' + P(x)y' + Q(x)y = g(x)
\]
\[
\frac{x^2 y'' - 2x y' + 2y}{x^2} = \frac{x^3 \ln x}{x^2}
\]
\[
y'' - \frac{2}{x}y' + \frac{2}{x^2}y = x \ln x
\]
\[
\boxed{g(x) = x \ln x}
\]

\vspace{3.5mm}

\textbf{\large 4. Calcular el Wronskiano}
\[
W(y_1, y_2)(x) = \begin{vmatrix} y_1 & y_2 \\ y_1' & y_2' \end{vmatrix} = \begin{vmatrix} x & x^2 \\ 1 & 2x \end{vmatrix}
\]
\[
W(x) = x(2x) - 1(x^2) = 2x^2 - x^2
\]
\[
\boxed{W(x) = x^2}
\]

\vspace{3.5mm}

\textbf{\large 5. Determinar los parámetros $u_1(x)$ y $u_2(x)$}
\[
\boxed{u_1'(x) = -\frac{y_2(x) g(x)}{W(x)}}
\]
\[
u_1'(x) = -\frac{x^2 (x \ln x)}{x^2} = -x \ln x
\]
\[
u_1(x) = -\int x \ln x\, dx = -\left(\frac{x^2}{2}\ln x - \int \frac{x^2}{2}\frac{1}{x}\, dx\right)
\]
\[
\boxed{u_1(x) = -\frac{1}{2}x^2 \ln x + \frac{1}{4}x^2}
\]
\[
\boxed{u_2'(x) = \frac{y_1(x) g(x)}{W(x)}}
\]
\[
u_2'(x) = \frac{x (x \ln x)}{x^2} = \ln x
\]
\[
u_2(x) = \int \ln x\, dx = x \ln x - \int x\left(\frac{1}{x}\right) dx
\]
\[
\boxed{u_2(x) = x \ln x - x}
\]

\vspace{3.5mm}

\textbf{\large 6. Construir la solución particular $y_p(x)$}
\[
\boxed{y_p(x) = u_1(x) y_1(x) + u_2(x) y_2(x)}
\]
\[
y_p(x) = \left(-\frac{1}{2}x^2 \ln x + \frac{1}{4}x^2\right) x + (x \ln x - x) x^2
\]
\[
y_p(x) = -\frac{1}{2}x^3 \ln x + \frac{1}{4}x^3 + x^3 \ln x - x^3
\]
\[
y_p(x) = \left(1 - \frac{1}{2}\right)x^3 \ln x + \left(\frac{1}{4} - 1\right)x^3
\]
\[
\boxed{y_p(x) = \frac{1}{2}x^3 \ln x - \frac{3}{4}x^3}
\]

\vspace{3.5mm}

\textbf{\large 7. Escribir la solución general}
\[
\boxed{y(x) = y_h(x) + y_p(x)}
\]
\[
\boxed{y(x) = c_1 x + c_2 x^2 + \frac{1}{2}x^3 \ln x - \frac{3}{4}x^3, \qquad x > 0}
\]

\begin{center}
\begin{tikzpicture}[xscale=1.1, yscale=0.42]
    % Ejes coordenados completos
    \draw[->, thick] (-0.8,0) -- (5.8,0) node[right] {$x$};
    \draw[->, thick] (0,-7.5) -- (0,8.5) node[above] {$y(x)$};
    
    % Marcas numéricas en los ejes
    \foreach \x in {1, 2, 3, 4, 5}
        \draw (\x, 0.15) -- (\x, -0.15) node[below, font=\scriptsize] {$\x$};
    \foreach \y in {-6, -4, -2, 2, 4, 6}
        \draw (0.1, \y) -- (-0.1, \y) node[left, font=\scriptsize] {$\y$};
        
    % Solución particular pura yp(x) (con c1=0, c2=0)
    \draw[very thick, deepteal, domain=0.08:5.1, samples=22] 
        plot (\x, {0.5*\x*\x*\x*ln(\x) - 0.75*\x*\x*\x});
        
    % Punto límite abierto en (0,0) ya que x > 0
    \draw[deepteal, thick, fill=white] (0,0) circle (2.2pt);
    
    % Punto mínimo local en x = e^{7/6} ≈ 3.21, y ≈ -5.52
    \filldraw[unmsmblue] (3.21, -5.52) circle (2.2pt) 
        node[below, font=\footnotesize] {$(3.21, -5.52)$};
        
    % Raíz / Cruce con el eje x en x = e^{1.5} ≈ 4.48
    \filldraw[deepteal] (4.48, 0) circle (2.2pt) 
        node[above left, font=\footnotesize] {$(4.48, 0)$};
        
    % Etiqueta de la curva
    \node[right, font=\footnotesize, deepteal] at (4.6, 5.5) {$y_p(x) = \frac{1}{2}x^3 \ln x - \frac{3}{4}x^3$};
\end{tikzpicture}
\end{center}

\end{solucion}

\vspace{3mm}

\begin{problema}[Problema 1.4]
Considere la ecuación diferencial no lineal de Riccati:
\[
\frac{dy}{dx} = y^2 - 2xy + (x^2 + 1)
\]
Verifique que $y_1(x) = x$ es una solución particular y determine la solución general $y(x)$.
\end{problema}

\begin{solucion}

\textbf{\large 1. Identificar la ecuación de Riccati}
\[
\boxed{y' = q_2(x)y^2 + q_1(x)y + q_0(x)}
\]
\[
y' = y^2 - 2xy + (x^2 + 1)
\]
\[
q_2(x) = 1, \qquad q_1(x) = -2x, \qquad q_0(x) = x^2 + 1
\]

\vspace{3.5mm}

\textbf{\large 2. Verificar la solución particular $y_1(x) = x$}
\[
y_1(x) = x \implies y_1'(x) = 1
\]
\[
y_1^2 - 2x y_1 + (x^2 + 1) = x^2 - 2x(x) + (x^2 + 1)
\]
\[
x^2 - 2x^2 + x^2 + 1 = 1
\]
\[
\boxed{y_1'(x) = 1 = y_1^2 - 2x y_1 + (x^2 + 1)}
\]

\vspace{3.5mm}

\textbf{\large 3. Aplicar el cambio de variable linealizante}
\[
\boxed{y(x) = y_1(x) + \frac{1}{u(x)} = x + \frac{1}{u(x)}}
\]
\[
y' = \frac{d}{dx}\left[x + u^{-1}\right]
\]
\[
\boxed{y' = 1 - \frac{u'}{u^2}}
\]

\vspace{3.5mm}

\textbf{\large 4. Sustituir en la ecuación diferencial}
\[
1 - \frac{u'}{u^2} = \left(x + \frac{1}{u}\right)^2 - 2x\left(x + \frac{1}{u}\right) + (x^2 + 1)
\]
\[
1 - \frac{u'}{u^2} = x^2 + \frac{2x}{u} + \frac{1}{u^2} - 2x^2 - \frac{2x}{u} + x^2 + 1
\]
\[
1 - \frac{u'}{u^2} = (x^2 - 2x^2 + x^2) + \left(\frac{2x}{u} - \frac{2x}{u}\right) + \frac{1}{u^2} + 1
\]
\[
1 - \frac{u'}{u^2} = 1 + \frac{1}{u^2}
\]

\vspace{3.5mm}

\textbf{\large 5. Resolver la ecuación diferencial para $u(x)$}
\[
-\frac{u'}{u^2} = \frac{1}{u^2}
\]
\[
\boxed{u' = -1}
\]
\[
\int du = \int (-1)\, dx
\]
\[
\boxed{u(x) = C - x}
\]

\vspace{3.5mm}

\textbf{\large 6. Construir la solución general $y(x)$}
\[
y(x) = x + \frac{1}{u(x)}
\]
\[
\boxed{y(x) = x + \frac{1}{C - x}}
\]
\[
y(x) = \frac{x(C - x) + 1}{C - x}
\]
\[
\boxed{y(x) = \frac{Cx - x^2 + 1}{C - x}, \qquad C \in \mathbb{R}}
\]

\begin{center}
\begin{tikzpicture}[xscale=0.9, yscale=0.38]
    % Ejes coordenados
    \draw[->, thick] (-0.8,0) -- (11.5,0) node[right] {$x$};
    \draw[->, thick] (0,-8.5) -- (0,10.8) node[above] {$y(x)$};
    
    % Marcas en los ejes
    \foreach \x in {2, 4, 6, 8, 10}
        \draw (\x, 0.2) -- (\x, -0.2) node[below, font=\scriptsize] {$\x$};
    \foreach \y in {-8, -6, -4, -2, 2, 4, 6, 8, 10}
        \draw (0.15, \y) -- (-0.15, \y) node[left, font=\scriptsize] {$\y$};
        
    % Asíntota vertical en x = 1 (con C = 1)
    \draw[dashed, gray!70, thick] (1, -8.2) -- (1, 10.2) 
        node[pos=0.96, left, font=\scriptsize, gray!90!black] {$x = 1$};
        
    % Solución particular conocida y1(x) = x para x > 0 (línea verde)
    \draw[very thick, green!60!black] (0.08, 0.08) -- (10.5, 10.5);
    \draw[green!60!black, thick, fill=white] (0, 0) circle (2.2pt);
    \node[above, font=\footnotesize, green!60!black] at (9.2, 9.4) {$y_1(x) = x$};
    
    % Rama 1 de la solución general con C = 1 (0 < x < 1)
    \draw[very thick, unmsmblue, domain=0.04:0.89, samples=16] 
        plot (\x, {\x + 1 / (1 - \x)});
    \draw[unmsmblue, thick, fill=white] (0, 1) circle (2.2pt) 
        node[left, font=\scriptsize, unmsmblue] {$(0, 1)$};
        
    % Rama 2 de la solución general con C = 1 (x > 1)
    \draw[very thick, unmsmblue, domain=1.11:10.5, samples=20] 
        plot (\x, {\x + 1 / (1 - \x)});
    \filldraw[unmsmblue] (2, 1) circle (2.2pt) 
        node[below right, font=\footnotesize] {$(2, 1)$};
    \node[below, font=\footnotesize, unmsmblue] at (5.5, 4.4) {$y(x) = x + \frac{1}{1-x}$};
\end{tikzpicture}
\end{center}

\end{solucion}

\vspace{3mm}

\begin{problema}[Problema 1.5]
El movimiento de un péndulo simple no lineal sin amortiguamiento se rige por:
\[
\ddot{\theta} + \omega_0^2 \sin\theta = 0
\]
Obtenga la primera integral de la energía, determine la ecuación de la separatriz en el plano $(\theta, \dot{\theta})$ y exprese el periodo exacto de oscilación para una amplitud angular $\theta_0 \in (0, \pi)$.
\end{problema}


\begin{solucion}

\textbf{\large 1. Obtener la primera integral de la energía}
\[
\ddot{\theta} + \omega_0^2 \sin\theta = 0
\]
\[
\dot{\theta}\ddot{\theta} + \omega_0^2 \sin\theta\,\dot{\theta} = 0
\]
\[
\frac{d}{dt}\left[ \frac{1}{2}\dot{\theta}^2 - \omega_0^2 \cos\theta \right] = 0
\]
\[
\boxed{E = \frac{1}{2}\dot{\theta}^2 - \omega_0^2 \cos\theta = \text{cte}}
\]
\[
\dot{\theta}(\theta_0) = 0 \implies E = -\omega_0^2 \cos\theta_0
\]
\[
\frac{1}{2}\dot{\theta}^2 - \omega_0^2 \cos\theta = -\omega_0^2 \cos\theta_0
\]
\[
\frac{1}{2}\dot{\theta}^2 = \omega_0^2(\cos\theta - \cos\theta_0)
\]
\[
\cos\theta - \cos\theta_0 = 2\left(\sin^2\frac{\theta_0}{2} - \sin^2\frac{\theta}{2}\right)
\]
\[
\boxed{\dot{\theta} = \pm 2\omega_0 \sqrt{\sin^2\left(\frac{\theta_0}{2}\right) - \sin^2\left(\frac{\theta}{2}\right)}}
\]

\vspace{3.5mm}

\textbf{\large 2. Determinar la ecuación de la separatriz}
\[
(\theta, \dot{\theta}) = (\pm\pi, 0)
\]
\[
E_{\text{sep}} = \frac{1}{2}(0)^2 - \omega_0^2 \cos(\pi) = \omega_0^2
\]
\[
\frac{1}{2}\dot{\theta}^2 - \omega_0^2 \cos\theta = \omega_0^2
\]
\[
\dot{\theta}^2 = 2\omega_0^2(1 + \cos\theta) = 4\omega_0^2 \cos^2\left(\frac{\theta}{2}\right)
\]
\[
\boxed{\dot{\theta}_{\text{sep}}(\theta) = \pm 2\omega_0 \cos\left(\frac{\theta}{2}\right)}
\]

\vspace{3.5mm}

\textbf{\large 3. Plantear la integral del periodo exacto}
\[
\frac{d\theta}{dt} = 2\omega_0 \sqrt{\sin^2\left(\frac{\theta_0}{2}\right) - \sin^2\left(\frac{\theta}{2}\right)}
\]
\[
dt = \frac{d\theta}{2\omega_0 \sqrt{\sin^2\left(\frac{\theta_0}{2}\right) - \sin^2\left(\frac{\theta}{2}\right)}}
\]
\[
T = 4 \int_0^{\theta_0} dt
\]
\[
\boxed{T = 4 \int_0^{\theta_0} \frac{d\theta}{2\omega_0 \sqrt{\sin^2\left(\frac{\theta_0}{2}\right) - \sin^2\left(\frac{\theta}{2}\right)}}}
\]

\vspace{3.5mm}

\textbf{\large 4. Aplicar la transformación de Jacobi}
\[
k = \sin\left(\frac{\theta_0}{2}\right)
\]
\[
\sin\left(\frac{\theta}{2}\right) = k \sin\phi
\]
\[
\frac{1}{2}\cos\left(\frac{\theta}{2}\right) d\theta = k \cos\phi\, d\phi
\]
\[
d\theta = \frac{2k\cos\phi\, d\phi}{\sqrt{1 - k^2\sin^2\phi}}
\]
\[
\theta = 0 \implies \phi = 0
\]
\[
\theta = \theta_0 \implies \phi = \frac{\pi}{2}
\]
\[
\sqrt{\sin^2\left(\frac{\theta_0}{2}\right) - \sin^2\left(\frac{\theta}{2}\right)} = k\cos\phi
\]

\vspace{3.5mm}

\textbf{\large 5. Expresar el periodo mediante la integral elíptica $K(k)$}
\[
T = \frac{4}{2\omega_0} \int_0^{\pi/2} \frac{1}{k\cos\phi} \left( \frac{2k\cos\phi\, d\phi}{\sqrt{1 - k^2\sin^2\phi}} \right)
\]
\[
T = \frac{4}{\omega_0} \int_0^{\pi/2} \frac{d\phi}{\sqrt{1 - k^2\sin^2\phi}}
\]
\[
K(k) = \int_0^{\pi/2} \frac{d\phi}{\sqrt{1 - k^2\sin^2\phi}}
\]
\[
\boxed{T = \frac{4}{\omega_0} K(k), \qquad k = \sin\left(\frac{\theta_0}{2}\right)}
\]
\[
\lim_{k \to 0} K(k) = \frac{\pi}{2} \implies \boxed{T_0 = \frac{2\pi}{\omega_0}}
\]

\begin{center}
\begin{tikzpicture}[xscale=0.7, yscale=0.72]
    % Leyenda superior tal como en la figura
    \node[above, font=\footnotesize] at (0, 3.8) {
        \textcolor{cyan!75!blue}{\rule{5mm}{1.5pt} \textbf{Oscilaciones ($E < 2\omega_0^2$)}}\qquad
        \textcolor{red!75!black}{\rule{5mm}{2pt} \textbf{Separatriz ($E = 2\omega_0^2$)}}\qquad
        \textcolor{green!65!black}{\rule{5mm}{1.5pt} \textbf{Rotaciones ($E > 2\omega_0^2$)}}
    };

    % Ejes coordenados
    \draw[->, thick] (-10.3,0) -- (10.6,0) node[right, font=\footnotesize] {$\theta$};
    \draw[->, thick] (0,-3.7) -- (0,3.7) node[above, font=\footnotesize] {$\dot{\theta}/\omega_0$};

    % Marcas en el eje theta (-3pi a 3pi)
    \draw (-9.4248, 0.15) -- (-9.4248, -0.15) node[below, font=\scriptsize] {$-3\pi$};
    \draw (-6.2832, 0.15) -- (-6.2832, -0.15) node[below, font=\scriptsize] {$-2\pi$};
    \draw (-3.1416, 0.15) -- (-3.1416, -0.15) node[below, font=\scriptsize] {$-\pi$};
    \draw (0, 0.15) -- (0, -0.15) node[below left, font=\scriptsize] {$0$};
    \draw (3.1416, 0.15) -- (3.1416, -0.15) node[below, font=\scriptsize] {$\pi$};
    \draw (6.2832, 0.15) -- (6.2832, -0.15) node[below, font=\scriptsize] {$2\pi$};
    \draw (9.4248, 0.15) -- (9.4248, -0.15) node[below, font=\scriptsize] {$3\pi$};

    % Marcas en el eje dot(theta)/omega0 (-3 a 3)
    \foreach \y in {-3, -2, -1, 1, 2, 3} {
        \draw (0.2, \y) -- (-0.2, \y) node[left, font=\scriptsize] {$\y$};
    }

    % Curvas de Rotación continua (verde, superior e inferior)
    \draw[thick, green!65!black, domain=-9.4248:9.4248, samples=15] 
        plot (\x, {sqrt(4.25 + 2*cos(\x r))});
    \draw[thick, green!65!black, domain=-9.4248:9.4248, samples=15] 
        plot (\x, {sqrt(7.0 + 2*cos(\x r))});
    \draw[thick, green!65!black, domain=-9.4248:9.4248, samples=15] 
        plot (\x, {-sqrt(4.25 + 2*cos(\x r))});
    \draw[thick, green!65!black, domain=-9.4248:9.4248, samples=15] 
        plot (\x, {-sqrt(7.0 + 2*cos(\x r))});

    % Curvas de Oscilación / Libración (azul, 5 órbitas anidadas elípticas en cada centro)
    \foreach \tc in {-6.2832, 0, 6.2832} {
        % Puntos de equilibrio estables (centros)
        \filldraw[cyan!75!blue] (\tc, 0) circle (2.5pt);

        \foreach \rx/\ry in {0.6/0.45, 1.1/0.85, 1.6/1.22, 2.1/1.52, 2.6/1.80} {
            \draw[thick, cyan!75!blue] (\tc, 0) ellipse (\rx and \ry);
        }
    }

    % Separatriz (rojo, ramas superior e inferior periódicas)
    \foreach \tc in {-6.2832, 0, 6.2832} {
        \draw[very thick, red!75!black, domain=\tc-3.1416:\tc+3.1416, samples=14] 
            plot (\x, {2 * cos(0.5*(\x - \tc) r)});
        \draw[very thick, red!75!black, domain=\tc-3.1416:\tc+3.1416, samples=14] 
            plot (\x, {-2 * cos(0.5*(\x - \tc) r)});
    }

    % Puntos de equilibrio inestables (puntos silla) con marcador característico
    \foreach \xs in {-9.4248, -3.1416, 3.1416, 9.4248} {
        \draw[red!75!black, thick, fill=white] (\xs, 0) circle (2.8pt);
        \draw[red!75!black, thick] (\xs-0.08, -0.15) -- (\xs+0.08, 0.15);
        \draw[red!75!black, thick] (\xs-0.08, 0.15) -- (\xs+0.08, -0.15);
    }
\end{tikzpicture}
\end{center}

\end{solucion}

\vspace{5mm}
% =========================================================================
% TEMA 2: RETRATOS DE FASE Y ANÁLISIS CUALITATIVO EN EL PLANO
% =========================================================================
\section*{\color{unmsmblue} Tema 2: Retratos de Fase y Análisis Cualitativo en el Plano}

\begin{problema}[Problema 2.1]
Considere el sistema dinámico lineal en el plano:
\[
\dot{x} = 2x, \qquad \dot{y} = -3y
\]
Clasifique el origen $(0,0)$, halle la ecuación de las trayectorias en coordenadas cartesianas y trace el comportamiento cualitativo del retrato de fase.
\end{problema}


\begin{solucion}

\textbf{\large 1. Representación matricial del sistema}
\[
\begin{pmatrix} \dot{x} \\ \dot{y} \end{pmatrix} = \begin{pmatrix} 2 & 0 \\ 0 & -3 \end{pmatrix} \begin{pmatrix} x \\ y \end{pmatrix}
\]
\[
\boxed{\dot{\mathbf{x}} = A \mathbf{x}, \qquad A = \begin{pmatrix} 2 & 0 \\ 0 & -3 \end{pmatrix}}
\]

\vspace{3.5mm}

\textbf{\large 2. Cálculo de autovalores y autovectores}
\[
\det(A - \lambda I) = 0
\]
\[
\begin{vmatrix} 2 - \lambda & 0 \\ 0 & -3 - \lambda \end{vmatrix} = (2 - \lambda)(-3 - \lambda) = 0
\]
\[
\boxed{\lambda_1 = 2, \qquad \lambda_2 = -3}
\]
\[
(A - \lambda_1 I)\mathbf{v}_1 = \mathbf{0} \implies \begin{pmatrix} 0 & 0 \\ 0 & -5 \end{pmatrix} \begin{pmatrix} v_{11} \\ v_{12} \end{pmatrix} = \begin{pmatrix} 0 \\ 0 \end{pmatrix} \implies \boxed{\mathbf{v}_1 = \begin{pmatrix} 1 \\ 0 \end{pmatrix}}
\]
\[
(A - \lambda_2 I)\mathbf{v}_2 = \mathbf{0} \implies \begin{pmatrix} 5 & 0 \\ 0 & 0 \end{pmatrix} \begin{pmatrix} v_{21} \\ v_{22} \end{pmatrix} = \begin{pmatrix} 0 \\ 0 \end{pmatrix} \implies \boxed{\mathbf{v}_2 = \begin{pmatrix} 0 \\ 1 \end{pmatrix}}
\]

\vspace{3.5mm}

\textbf{\large 3. Clasificación del punto de equilibrio $(0,0)$}
\[
\det(A) = (2)(-3) = -6 < 0
\]
\[
\lambda_2 = -3 < 0 < 2 = \lambda_1
\]
\[
\boxed{(0,0) \text{ es un Punto Silla (Punto de Ensilladura) Inestable}}
\]

\vspace{3.5mm}

\textbf{\large 4. Obtención de las trayectorias cartesianas (Órbitas)}
\[
\dot{x} = 2x \implies x(t) = x_0 e^{2t}
\]
\[
\dot{y} = -3y \implies y(t) = y_0 e^{-3t}
\]
\[
\frac{dy}{dx} = \frac{\dot{y}}{\dot{x}} = \frac{-3y}{2x}
\]
\[
\frac{dy}{y} = -\frac{3}{2}\frac{dx}{x}
\]
\[
\int \frac{dy}{y} = -\frac{3}{2}\int \frac{dx}{x}
\]
\[
\ln|y| = -\frac{3}{2}\ln|x| + C_1
\]
\[
\ln|y| + \ln|x|^{3/2} = C_1 \implies \ln\left(|y|\,|x|^{3/2}\right) = C_1
\]
\[
|y|\,|x|^{3/2} = K
\]
\[
\boxed{y^2 x^3 = C, \qquad C \ge 0}
\]
\[
\boxed{y(x) = \pm \frac{C}{\sqrt{x^3}} = \pm C x^{-3/2}}
\]

\vspace{3.5mm}

\textbf{\large 5. Caracterización de las variedades invariantes $W^s$ y $W^u$}
\[
W^s(0,0) = \mathrm{gen}\{\mathbf{v}_2\} = \left\{ (x, y) \in \mathbb{R}^2 : x = 0 \right\} \quad (\text{Eje } y)
\]
\[
\lim_{t \to +\infty} \mathbf{x}(t) = (0, 0), \qquad \forall (0, y_0) \in W^s
\]
\[
\boxed{W^s(0,0) = \left\{ (0, y) \in \mathbb{R}^2 \right\} \quad (\text{Variedad Estable})}
\]
\[
W^u(0,0) = \mathrm{gen}\{\mathbf{v}_1\} = \left\{ (x, y) \in \mathbb{R}^2 : y = 0 \right\} \quad (\text{Eje } x)
\]
\[
\lim_{t \to -\infty} \mathbf{x}(t) = (0, 0), \qquad \forall (x_0, 0) \in W^u
\]
\[
\boxed{W^u(0,0) = \left\{ (x, 0) \in \mathbb{R}^2 \right\} \quad (\text{Variedad Inestable})}
\]

\begin{center}
\begin{tikzpicture}[scale=1.1]
    % Variedad inestable (y = 0, eje horizontal rojo)
    \draw[very thick, red!85!black] (-3.8, 0) -- (3.8, 0) node[right, black, font=\footnotesize] {$x$};
    % Flechas en variedad inestable saliendo de (0,0)
    \draw[-{Latex[length=2.8mm]}, very thick, red!85!black] (0, 0) -- (2.0, 0);
    \draw[-{Latex[length=2.8mm]}, very thick, red!85!black] (0, 0) -- (-2.0, 0);
    \node[below, font=\scriptsize, red!85!black] at (2.4, -0.05) {Variedad inestable $(y = 0)$};

    % Variedad estable (x = 0, eje vertical azul)
    \draw[very thick, cyan!80!blue] (0, -3.2) -- (0, 3.2) node[above, black, font=\footnotesize] {$y$};
    % Flechas en variedad estable entrando hacia (0,0)
    \draw[-{Latex[length=2.8mm]}, very thick, cyan!80!blue] (0, 2.8) -- (0, 1.4);
    \draw[-{Latex[length=2.8mm]}, very thick, cyan!80!blue] (0, -2.8) -- (0, -1.4);
    \node[right, font=\scriptsize, cyan!80!blue] at (0.1, -2.9) {Variedad estable $(x = 0)$};

    % Trayectorias de fase x^3 y^2 = C (3 en cada cuadrante, líneas continuas)
    % Cuadrante I
    \draw[thick, gray!70!black, domain=0.28:3.6, samples=12, smooth] plot (\x, {0.32 / (\x * sqrt(\x))});
    \draw[thick, gray!70!black, domain=0.48:3.6, samples=12, smooth] plot (\x, {0.75 / (\x * sqrt(\x))});
    \draw[thick, gray!70!black, domain=0.75:3.6, samples=12, smooth] plot (\x, {1.55 / (\x * sqrt(\x))});
    % Flechas Cuadrante I
    \draw[-{Latex[length=2.2mm]}, gray!70!black] (0.5, 0.90) -- (0.8, 0.45);
    \draw[-{Latex[length=2.2mm]}, gray!70!black] (0.9, 1.0) -- (1.3, 0.55);
    \draw[-{Latex[length=2.2mm]}, gray!70!black] (1.3, 1.1) -- (1.8, 0.65);

    % Cuadrante II
    \draw[thick, gray!70!black, domain=0.28:3.6, samples=12, smooth] plot (-\x, {0.32 / (\x * sqrt(\x))});
    \draw[thick, gray!70!black, domain=0.48:3.6, samples=12, smooth] plot (-\x, {0.75 / (\x * sqrt(\x))});
    \draw[thick, gray!70!black, domain=0.75:3.6, samples=12, smooth] plot (-\x, {1.55 / (\x * sqrt(\x))});
    % Flechas Cuadrante II
    \draw[-{Latex[length=2.2mm]}, gray!70!black] (-0.5, 0.90) -- (-0.8, 0.45);
    \draw[-{Latex[length=2.2mm]}, gray!70!black] (-0.9, 1.0) -- (-1.3, 0.55);
    \draw[-{Latex[length=2.2mm]}, gray!70!black] (-1.3, 1.1) -- (-1.8, 0.65);

    % Cuadrante III
    \draw[thick, gray!70!black, domain=0.28:3.6, samples=12, smooth] plot (-\x, {-0.32 / (\x * sqrt(\x))});
    \draw[thick, gray!70!black, domain=0.48:3.6, samples=12, smooth] plot (-\x, {-0.75 / (\x * sqrt(\x))});
    \draw[thick, gray!70!black, domain=0.75:3.6, samples=12, smooth] plot (-\x, {-1.55 / (\x * sqrt(\x))});
    % Flechas Cuadrante III
    \draw[-{Latex[length=2.2mm]}, gray!70!black] (-0.8, -0.45) -- (-1.3, -0.22);
    \draw[-{Latex[length=2.2mm]}, gray!70!black] (-1.1, -0.7) -- (-1.7, -0.35);
    \draw[-{Latex[length=2.2mm]}, gray!70!black] (-1.5, -0.95) -- (-2.1, -0.5);

    % Cuadrante IV
    \draw[thick, gray!70!black, domain=0.28:3.6, samples=12, smooth] plot (\x, {-0.32 / (\x * sqrt(\x))});
    \draw[thick, gray!70!black, domain=0.48:3.6, samples=12, smooth] plot (\x, {-0.75 / (\x * sqrt(\x))});
    \draw[thick, gray!70!black, domain=0.75:3.6, samples=12, smooth] plot (\x, {-1.55 / (\x * sqrt(\x))});
    % Flechas Cuadrante IV
    \draw[-{Latex[length=2.2mm]}, gray!70!black] (0.8, -0.45) -- (1.3, -0.22);
    \draw[-{Latex[length=2.2mm]}, gray!70!black] (1.1, -0.7) -- (1.7, -0.35);
    \draw[-{Latex[length=2.2mm]}, gray!70!black] (1.5, -0.95) -- (2.1, -0.5);

    % Etiqueta x^3 y^2 = C
    \node[font=\footnotesize\bfseries, gray!85!black] at (2.4, 2.4) {$x^3 y^2 = C$};

    % Punto de equilibrio silla en el origen
    \draw[black, very thick, fill=white] (0, 0) circle (3pt);
    \node[below right, font=\footnotesize\bfseries] at (0.1, -0.05) {$(0, 0)$ silla};
\end{tikzpicture}
\end{center}

\end{solucion}

\vspace{3mm}

\begin{problema}[Problema 2.2]
Clasifique el equilibrio $(0,0)$ del oscilador lineal según el parámetro $\gamma \ge 0$:
\[
\dot{x} = y, \qquad \dot{y} = -4x - \gamma y
\]
\end{problema}


\begin{solucion}

\textbf{\large 1. Representación matricial del sistema}
\[
\begin{pmatrix} \dot{x} \\ \dot{y} \end{pmatrix} = \begin{pmatrix} 0 & 1 \\ -4 & -\gamma \end{pmatrix} \begin{pmatrix} x \\ y \end{pmatrix}
\]
\[
\boxed{\dot{\mathbf{x}} = A \mathbf{x}, \qquad A = \begin{pmatrix} 0 & 1 \\ -4 & -\gamma \end{pmatrix}}
\]

\vspace{3.5mm}

\textbf{\large 2. Cálculo de la traza, determinante y ecuación característica}
\[
\tau = \mathrm{Tr}(A) = 0 + (-\gamma) = -\gamma
\]
\[
\Delta = \det(A) = (0)(-\gamma) - (1)(-4) = 4
\]
\[
\det(A - \lambda I) = \lambda^2 - \tau \lambda + \Delta = 0
\]
\[
\boxed{\lambda^2 + \gamma \lambda + 4 = 0}
\]

\vspace{3.5mm}

\textbf{\large 3. Análisis del discriminante y autovalores}
\[
D = \tau^2 - 4\Delta = \gamma^2 - 4(4)
\]
\[
\boxed{D = \gamma^2 - 16}
\]
\[
\boxed{\lambda_{1, 2} = \frac{-\gamma \pm \sqrt{\gamma^2 - 16}}{2}}
\]

\vspace{3.5mm}

\textbf{\large 4. Clasificación según el parámetro $\gamma \ge 0$}
\[
\gamma = 0 \implies \tau = 0, \quad D = -16 < 0 \implies \lambda_{1,2} = \pm 2i
\]
\[
\boxed{\gamma = 0 \implies \text{Centro Neutramente Estable (Órbitas Elípticas Cerradas)}}
\]
\[
0 < \gamma < 4 \implies \tau < 0, \quad D < 0 \implies \lambda_{1,2} = -\frac{\gamma}{2} \pm i \frac{\sqrt{16-\gamma^2}}{2}
\]
\[
\boxed{0 < \gamma < 4 \implies \text{Foco Estable (Espiral Atractora Asintóticamente Estable)}}
\]
\[
\gamma = 4 \implies \tau = -4 < 0, \quad D = 0 \implies \lambda_1 = \lambda_2 = -2
\]
\[
\boxed{\gamma = 4 \implies \text{Nodo Impropio Degenerado Estable (Amortiguamiento Crítico)}}
\]
\[
\gamma > 4 \implies \tau < 0, \quad D > 0 \implies \lambda_2 < \lambda_1 < 0
\]
\[
\boxed{\gamma > 4 \implies \text{Nodo Estable (Sobreamortiguado, Dos Raíces Reales Negativas)}}
\]

\begin{center}
\begin{tikzpicture}[scale=0.95]
    % Ejes coordenados
    \draw[->, thick] (-3.2,0) -- (3.2,0) node[right] {$x$};
    \draw[->, thick] (0,-3.2) -- (0,3.2) node[above] {$y$};
    
    % Órbita elíptica cerrada para gamma = 0 (Centro, línea continua verde)
    \draw[thick, green!60!black] (0,0) ellipse (1.2cm and 2.4cm);
    \draw[-{Latex[length=2.2mm]}, green!60!black] (1.2, 0.05) -- (1.2, -0.15);
    \draw[-{Latex[length=2.2mm]}, green!60!black] (-1.2, -0.05) -- (-1.2, 0.15);
    \node[right, font=\footnotesize, green!60!black] at (1.2, 1.8) {$\gamma = 0 \text{ (Centro)}$};
    
    % Espirales atractoras para 0 < gamma < 4 (Foco estable, sentido horario)
    \draw[domain=0:20, samples=20, smooth, very thick, deepteal, -{Latex[length=2.5mm]}] 
        plot ({2.6*exp(-0.14*\x)*cos(-\x r)}, {-2.6*exp(-0.14*\x)*sin(-\x r)});
    \draw[domain=0:20, samples=20, smooth, thick, unmsmblue, -{Latex[length=2.5mm]}] 
        plot ({-2.2*exp(-0.14*\x)*cos(-\x r)}, {2.2*exp(-0.14*\x)*sin(-\x r)});
    \node[above left, font=\footnotesize, deepteal] at (-1.5, 2.2) {$0 < \gamma < 4 \text{ (Foco)}$};
    
    % Punto de equilibrio en el origen
    \filldraw[black] (0,0) circle (2.2pt) node[below right, font=\footnotesize] {$(0,0)$};
\end{tikzpicture}
\end{center}

\end{solucion}

\vspace{3mm}

\begin{problema}[Problema 2.3]
Analice los puntos de equilibrio, la linealización y la estabilidad del sistema:
\[
\dot{x} = x(3 - x - 2y), \qquad \dot{y} = y(2 - x - y)
\]
en el cuadrante biológico $x \ge 0, y \ge 0$.
\end{problema}

\begin{solucion}

\textbf{\large 1. Determinar los puntos de equilibrio en el cuadrante biológico}
\[
\dot{x} = x(3 - x - 2y) = 0 \iff x = 0 \quad \lor \quad x + 2y = 3
\]
\[
\dot{y} = y(2 - x - y) = 0 \iff y = 0 \quad \lor \quad x + y = 2
\]

\textbf{Punto de extinción total:}
\[
x = 0, \quad y = 0 \implies \boxed{P_1 = (0, 0)}
\]

\textbf{Punto de exclusión de la especie $x$ (frontera vertical):}
\[
x = 0, \quad 2 - 0 - y = 0 \implies \boxed{P_2 = (0, 2)}
\]

\textbf{Punto de exclusión de la especie $y$ (frontera horizontal):}
\[
y = 0, \quad 3 - x - 2(0) = 0 \implies \boxed{P_3 = (3, 0)}
\]

\textbf{Punto de coexistencia interior:}
\[
\begin{cases}
x + 2y = 3 \\
x + y = 2
\end{cases}
\implies (x + 2y) - (x + y) = 3 - 2 \implies y = 1, \quad x = 2 - 1 = 1
\]
\[
\boxed{P_4 = (1, 1)}
\]

\vspace{3.5mm}

\textbf{\large 2. Calcular la matriz jacobiana general}
\[
f(x, y) = 3x - x^2 - 2xy \implies \frac{\partial f}{\partial x} = 3 - 2x - 2y, \quad \frac{\partial f}{\partial y} = -2x
\]
\[
g(x, y) = 2y - xy - y^2 \implies \frac{\partial g}{\partial x} = -y, \quad \frac{\partial g}{\partial y} = 2 - x - 2y
\]
\[
\boxed{J(x, y) = \begin{pmatrix} 3 - 2x - 2y & -2x \\ -y & 2 - x - 2y \end{pmatrix}}
\]

\vspace{3.5mm}

\textbf{\large 3. Clasificación y estabilidad de cada punto de equilibrio}

\textbf{Análisis en $P_1(0, 0)$:}
\[
J(0, 0) = \begin{pmatrix} 3 & 0 \\ 0 & 2 \end{pmatrix}
\]
\[
\det(J - \lambda I) = (3 - \lambda)(2 - \lambda) = 0 \implies \lambda_1 = 3, \quad \lambda_2 = 2
\]
\[
\boxed{\lambda_1 = 3 > 0, \quad \lambda_2 = 2 > 0 \implies \text{Nodo Inestable (Fuente)}}
\]

\textbf{Análisis en $P_2(0, 2)$:}
\[
J(0, 2) = \begin{pmatrix} 3 - 4 & 0 \\ -2 & 2 - 4 \end{pmatrix} = \begin{pmatrix} -1 & 0 \\ -2 & -2 \end{pmatrix}
\]
\[
\det(J - \lambda I) = (-1 - \lambda)(-2 - \lambda) = 0 \implies \lambda_1 = -1, \quad \lambda_2 = -2
\]
\[
\boxed{\lambda_1 = -1 < 0, \quad \lambda_2 = -2 < 0 \implies \text{Nodo Asintóticamente Estable (Sumidero)}}
\]

\textbf{Análisis en $P_3(3, 0)$:}
\[
J(3, 0) = \begin{pmatrix} 3 - 6 & -6 \\ 0 & 2 - 3 \end{pmatrix} = \begin{pmatrix} -3 & -6 \\ 0 & -1 \end{pmatrix}
\]
\[
\det(J - \lambda I) = (-3 - \lambda)(-1 - \lambda) = 0 \implies \lambda_1 = -3, \quad \lambda_2 = -1
\]
\[
\boxed{\lambda_1 = -3 < 0, \quad \lambda_2 = -1 < 0 \implies \text{Nodo Asintóticamente Estable (Sumidero)}}
\]

\textbf{Análisis en $P_4(1, 1)$:}
\[
J(1, 1) = \begin{pmatrix} 3 - 2 - 2 & -2 \\ -1 & 2 - 1 - 2 \end{pmatrix} = \begin{pmatrix} -1 & -2 \\ -1 & -1 \end{pmatrix}
\]
\[
\mathrm{Tr}(J) = -2, \qquad \det(J) = (-1)(-1) - (-2)(-1) = 1 - 2 = -1 < 0
\]
\[
\lambda^2 - \mathrm{Tr}(J)\lambda + \det(J) = 0 \implies \lambda^2 + 2\lambda - 1 = 0
\]
\[
\lambda = \frac{-2 \pm \sqrt{4 - 4(1)(-1)}}{2} = \frac{-2 \pm 2\sqrt{2}}{2} = -1 \pm \sqrt{2}
\]
\[
\boxed{\lambda_1 = -1 - \sqrt{2} \approx -2.414 < 0, \quad \lambda_2 = -1 + \sqrt{2} \approx +0.414 > 0 \implies \text{Punto Silla (Inestable)}}
\]

\textbf{Direcciones propias en $P_4(1, 1)$:}
\[
(J - \lambda I)\mathbf{v} = \mathbf{0} \implies \begin{pmatrix} -1 - \lambda & -2 \\ -1 & -1 - \lambda \end{pmatrix}\begin{pmatrix} v_1 \\ v_2 \end{pmatrix} = \begin{pmatrix} 0 \\ 0 \end{pmatrix} \implies v_1 + (1+\lambda)v_2 = 0
\]
\[
\mathbf{v}^{(s)} = \begin{pmatrix} \sqrt{2} \\ 1 \end{pmatrix} \approx \begin{pmatrix} 1.414 \\ 1 \end{pmatrix} \quad (\text{Pendiente } m_s = \frac{1}{\sqrt{2}} \approx 0.707 > 0)
\]
\[
\mathbf{v}^{(u)} = \begin{pmatrix} -\sqrt{2} \\ 1 \end{pmatrix} \approx \begin{pmatrix} -1.414 \\ 1 \end{pmatrix} \quad (\text{Pendiente } m_u = -\frac{1}{\sqrt{2}} \approx -0.707 < 0)
\]

\vspace{3.5mm}

\textbf{\large 4. Identificar las nullclinas y regiones de flujo}
\[
\boxed{\mathcal{N}_x \, (\dot{x} = 0): \quad x = 0 \quad \lor \quad x + 2y = 3 \quad (\text{Flujo estrictamente vertical})}
\]
\[
\boxed{\mathcal{N}_y \, (\dot{y} = 0): \quad y = 0 \quad \lor \quad x + y = 2 \quad (\text{Flujo estrictamente horizontal})}
\]

\vspace{3.5mm}

\textbf{\large 5. Retrato de fase en el cuadrante biológico $x \ge 0, y \ge 0$}

\begin{center}
\begin{tikzpicture}[xscale=1.5, yscale=1.6]
    % Regiones de fondo sutiles para cuencas de atracción
    \fill[solgreen!45, opacity=0.6] (0, 0.42) to[out=28, in=215] (1, 1) to[out=35, in=205] (3.6, 2.3) -- (3.6, 2.8) -- (0, 2.8) -- cycle;
    \fill[probhead!55, opacity=0.6] (0, 0.42) to[out=28, in=215] (1, 1) to[out=35, in=205] (3.6, 2.3) -- (3.6, 0) -- (0, 0) -- cycle;

    % Ejes coordenados del cuadrante biológico
    \draw[->, thick] (-0.15, 0) -- (3.8, 0) node[right, font=\footnotesize] {$x$ (Especie 1)};
    \draw[->, thick] (0, -0.15) -- (0, 2.7) node[above, font=\footnotesize] {$y$ (Especie 2)};

    % Graduaciones en los ejes
    \foreach \x in {1, 2, 3}
        \draw (\x, 0.06) -- (\x, -0.06) node[below, font=\scriptsize] {$\x$};
    \foreach \y in {1, 2}
        \draw (0.04, \y) -- (-0.04, \y) node[left, font=\scriptsize] {$\y$};

    % Nullclina x: x + 2y = 3 (color deepteal)
    \draw[very thick, deepteal] (0, 1.5) -- (3.3, -0.15);
    \node[deepteal, font=\scriptsize\bfseries, above right] at (2.2, 0.42) {$\dot{x}=0: x+2y=3$};

    % Vectores de flujo vertical sobre la nullclina x
    \draw[-{Latex[length=1.8mm]}, deepteal, thick] (0.6, 1.2) -- (0.6, 0.95);
    \draw[-{Latex[length=1.8mm]}, deepteal, thick] (1.6, 0.7) -- (1.6, 0.9);
    \draw[-{Latex[length=1.8mm]}, deepteal, thick] (2.4, 0.3) -- (2.4, 0.5);

    % Nullclina y: x + y = 2 (color unmsmblue)
    \draw[very thick, unmsmblue] (0, 2.0) -- (2.4, -0.4);
    \node[unmsmblue, font=\scriptsize\bfseries, above right] at (0.35, 1.72) {$\dot{y}=0: x+y=2$};

    % Vectores de flujo horizontal sobre la nullclina y
    \draw[-{Latex[length=1.8mm]}, unmsmblue, thick] (0.5, 1.5) -- (0.28, 1.5);
    \draw[-{Latex[length=1.8mm]}, unmsmblue, thick] (1.4, 0.6) -- (1.65, 0.6);

    % Separatriz de cuencas (Variedad estable del punto silla W^s(P4))
    \draw[line width=1.2pt, darkpurple, dashed] (0, 0.42) to[out=28, in=215] (1, 1) to[out=35, in=205] (3.6, 2.3);
    \draw[-{Latex[length=2.5mm]}, darkpurple, very thick] (0.45, 0.68) -- (0.60, 0.77);
    \draw[-{Latex[length=2.5mm]}, darkpurple, very thick] (2.6, 1.83) -- (2.35, 1.70);
    \node[darkpurple, font=\scriptsize\bfseries, above, rotate=28] at (2.45, 1.85) {Separatriz $W^s(P_4)$};

    % Variedad inestable saliendo de P4 hacia P2 y P3
    \draw[line width=1.1pt, red!75!black] (1, 1) to[out=145, in=-40] (0.12, 1.85);
    \draw[-{Latex[length=2.2mm]}, red!75!black, thick] (0.6, 1.45) -- (0.45, 1.62);
    \draw[line width=1.1pt, red!75!black] (1, 1) to[out=-35, in=145] (2.85, 0.12);
    \draw[-{Latex[length=2.2mm]}, red!75!black, thick] (1.8, 0.55) -- (2.1, 0.40);

    % Trayectorias en la Cuenca Superior (convergen a P2)
    \draw[thick, solborder, -{Latex[length=2.0mm]}] (0.6, 2.4) to[out=-75, in=15] (0.1, 2.05);
    \draw[thick, solborder, -{Latex[length=2.0mm]}] (1.8, 2.5) to[out=-120, in=25] (0.15, 2.12);
    \draw[thick, solborder, -{Latex[length=2.0mm]}] (0.2, 0.7) to[out=80, in=-95] (0.05, 1.85);

    % Trayectorias en la Cuenca Inferior (convergen a P3)
    \draw[thick, solborder, -{Latex[length=2.0mm]}] (3.4, 1.4) to[out=-135, in=45] (3.08, 0.15);
    \draw[thick, solborder, -{Latex[length=2.0mm]}] (2.6, 2.3) to[out=-115, in=75] (3.12, 0.25);
    \draw[thick, solborder, -{Latex[length=2.0mm]}] (1.1, 0.3) to[out=-10, in=165] (2.85, 0.05);

    % Texto identificador de cuencas
    \node[solborder, font=\scriptsize\bfseries] at (0.9, 2.35) {Cuenca de Atracción de $P_2$};
    \node[unmsmblue, font=\scriptsize\bfseries] at (2.9, 0.85) {Cuenca de Atracción de $P_3$};

    % Puntos de equilibrio con marcadores formales
    % P1 (0,0) - Fuente
    \draw[black, thick, fill=white] (0, 0) circle (2.5pt);
    \node[below left, font=\scriptsize\bfseries] at (0, 0) {$P_1(0,0)$};

    % P2 (0,2) - Sumidero
    \filldraw[solborder] (0, 2) circle (2.8pt);
    \node[left, font=\scriptsize\bfseries, solborder] at (-0.08, 2) {$P_2(0,2)$};

    % P3 (3,0) - Sumidero
    \filldraw[solborder] (3, 0) circle (2.8pt);
    \node[below, font=\scriptsize\bfseries, solborder] at (3, -0.08) {$P_3(3,0)$};

    % P4 (1,1) - Silla
    \draw[black, thick, fill=red!80!black] (1, 1) circle (2.8pt);
    \node[above right, font=\scriptsize\bfseries, red!80!black] at (1.02, 1.05) {$P_4(1,1)$ Silla};
\end{tikzpicture}
\end{center}

\end{solucion}

\begin{conclusion}
El punto de coexistencia es inestable. Las cuencas de atracción de los dos equilibrios estables $P_2(0,2)$ y $P_3(3,0)$ están separadas por la variedad estable del punto silla $P_4(1,1)$ (Principio de Exclusión Competitiva).
\end{conclusion}

\vspace{3mm}

\begin{problema}[Problema 2.4]
Considere el sistema dinámico conservativo:
\[
\dot{x} = y, \qquad \dot{y} = x - x^3
\]
Halle su función hamiltoniana $H(x,y)$, clasifique sus puntos fijos y determine la ecuación exacta de la órbita homoclínica (separatriz).
\end{problema}

\begin{solucion}

\textbf{\large 1. Determinar la función hamiltoniana $H(x, y)$}
\[
\dot{x} = \frac{\partial H}{\partial y} = y \implies H(x, y) = \frac{1}{2}y^2 + V(x)
\]
\[
\dot{y} = -\frac{\partial H}{\partial x} = -V'(x) = x - x^3 \implies V'(x) = x^3 - x
\]
\[
V(x) = \int (x^3 - x)\, dx = \frac{1}{4}x^4 - \frac{1}{2}x^2
\]
\[
\boxed{H(x, y) = \frac{1}{2}y^2 - \frac{1}{2}x^2 + \frac{1}{4}x^4}
\]

\textbf{Conservación de la energía del sistema:}
\[
\frac{dH}{dt} = \frac{\partial H}{\partial x}\dot{x} + \frac{\partial H}{\partial y}\dot{y} = (x^3 - x)y + y(x - x^3) \equiv 0 \implies \boxed{H(x(t), y(t)) = E = \text{cte}}
\]

\vspace{3.5mm}

\textbf{\large 2. Hallar y clasificar los puntos de equilibrio}
\[
\dot{x} = y = 0
\]
\[
\dot{y} = x(1 - x^2) = 0 \iff x = 0 \quad \lor \quad x = 1 \quad \lor \quad x = -1
\]
\[
\boxed{P_0 = (0, 0), \qquad P_+ = (1, 0), \qquad P_- = (-1, 0)}
\]

\textbf{Matriz jacobiana general:}
\[
J(x, y) = \begin{pmatrix} 0 & 1 \\ 1 - 3x^2 & 0 \end{pmatrix}, \qquad \mathrm{Tr}(J) \equiv 0
\]

\textbf{Análisis en el origen $P_0(0, 0)$:}
\[
J(0, 0) = \begin{pmatrix} 0 & 1 \\ 1 & 0 \end{pmatrix}
\]
\[
\det(J - \lambda I) = \lambda^2 - 1 = 0 \implies \lambda_1 = -1, \quad \lambda_2 = +1
\]
\[
\boxed{\lambda_{1,2} = \pm 1 \implies \text{Punto Silla (Inestable)}}
\]
\[
\boxed{H(0, 0) = 0}
\]

\textbf{Análisis en los puntos simétricos $P_{\pm}(\pm 1, 0)$:}
\[
J(\pm 1, 0) = \begin{pmatrix} 0 & 1 \\ 1 - 3 & 0 \end{pmatrix} = \begin{pmatrix} 0 & 1 \\ -2 & 0 \end{pmatrix}
\]
\[
\det(J - \lambda I) = \lambda^2 + 2 = 0 \implies \lambda_{1,2} = \pm i\sqrt{2}
\]
\[
V''(x) = 3x^2 - 1 \implies V''(\pm 1) = 2 > 0 \quad (\text{Mínimo local estricto de energía potencial})
\]
\[
\boxed{\lambda_{1,2} = \pm i\sqrt{2} \implies \text{Centros No Lineales (Estables en el sentido de Lyapunov)}}
\]
\[
\boxed{H(\pm 1, 0) = -\frac{1}{2}(1)^2 + \frac{1}{4}(1)^4 = -\frac{1}{4}}
\]

\vspace{3.5mm}

\textbf{\large 3. Obtener la ecuación analítica de la separatriz homoclínica}
\[
H(x, y) = H(0, 0) = 0
\]
\[
\frac{1}{2}y^2 - \frac{1}{2}x^2 + \frac{1}{4}x^4 = 0 \iff y^2 = x^2\left(1 - \frac{x^2}{2}\right)
\]
\[
1 - \frac{x^2}{2} \ge 0 \iff x^2 \le 2 \iff -\sqrt{2} \le x \le \sqrt{2}
\]
\[
\boxed{y_{\text{hom}}(x) = \pm x\sqrt{1 - \frac{x^2}{2}}, \qquad x \in [-\sqrt{2}, \, \sqrt{2}]}
\]

\textbf{Propiedades geométricas de los bucles homoclínicos:}
\[
x_{\text{vértice}} = \pm\sqrt{2} \approx \pm 1.414 \quad (y = 0)
\]
\[
\frac{d(y^2)}{dx} = 2x - 2x^3 = 0 \implies x = \pm 1 \implies y_{\text{máx}} = \pm\frac{1}{\sqrt{2}} \approx \pm 0.707
\]
\[
\boxed{\text{Vértices en } (\pm\sqrt{2}, 0), \qquad \text{Alturas extremas en } (\pm 1, \, \pm 1/\sqrt{2})}
\]


\textbf{\large 4. Retrato de fase hamiltoniano en el plano $(x, y)$}

\begin{center}
\begin{tikzpicture}[xscale=1.5, yscale=1.35]
    % Curvas de nivel abiertas exteriores (E = 0.25 > 0)
    \draw[thick, gray!60, domain=-1.52:1.52, samples=15, smooth] 
        plot (\x, {sqrt(max(0, 0.5 + \x*\x - 0.5*\x*\x*\x*\x))});
    \draw[thick, gray!60, domain=-1.52:1.52, samples=15, smooth] 
        plot (\x, {-sqrt(max(0, 0.5 + \x*\x - 0.5*\x*\x*\x*\x))});
    \draw[-{Latex[length=1.8mm]}, gray!80] (0.1, 0.71) -- (0.35, 0.71);
    \draw[-{Latex[length=1.8mm]}, gray!80] (-0.1, -0.71) -- (-0.35, -0.71);
    \node[right, font=\tiny, gray!80] at (1.55, 0.35) {$H > 0$};

    % Ejes coordenados
    \draw[->, thick] (-2.2, 0) -- (2.2, 0) node[right, font=\footnotesize] {$x$};
    \draw[->, thick] (0, -1.25) -- (0, 1.3) node[above, font=\footnotesize] {$y = \dot{x}$};

    % Graduaciones en los ejes
    \draw (1, 0.05) -- (1, -0.05) node[below, font=\scriptsize] {$1$};
    \draw (-1, 0.05) -- (-1, -0.05) node[below, font=\scriptsize] {$-1$};
    \draw (1.414, 0.05) -- (1.414, -0.05) node[below right, font=\tiny] {$\sqrt{2}$};
    \draw (-1.414, 0.05) -- (-1.414, -0.05) node[below left, font=\tiny] {$-\sqrt{2}$};
    \draw (0.04, 0.707) -- (-0.04, 0.707) node[left, font=\tiny] {$\frac{1}{\sqrt{2}}$};
    \draw (0.04, -0.707) -- (-0.04, -0.707) node[left, font=\tiny] {$-\frac{1}{\sqrt{2}}$};

    % Familias de órbitas cerradas periódicas interiores (Centros)
    % Centro derecho (1, 0)
    \draw[thick, cyan!75!blue] (1, 0) ellipse (0.28 and 0.20);
    \draw[thick, cyan!75!blue] (1, 0) ellipse (0.52 and 0.38);
    \draw[-{Latex[length=1.8mm]}, cyan!75!blue, thick] (1.0, 0.38) -- (1.15, 0.38);
    \draw[-{Latex[length=1.8mm]}, cyan!75!blue, thick] (1.0, -0.38) -- (0.85, -0.38);

    % Centro izquierdo (-1, 0)
    \draw[thick, cyan!75!blue] (-1, 0) ellipse (0.28 and 0.20);
    \draw[thick, cyan!75!blue] (-1, 0) ellipse (0.52 and 0.38);
    \draw[-{Latex[length=1.8mm]}, cyan!75!blue, thick] (-1.0, 0.38) -- (-0.85, 0.38);
    \draw[-{Latex[length=1.8mm]}, cyan!75!blue, thick] (-1.0, -0.38) -- (-1.15, -0.38);

    % Bucle homoclínico en forma de ocho (Separatriz H = 0)
    % Lóbulo derecho (x in [0, 1.414])
    \draw[very thick, red!80!black, domain=0:1.414, samples=15, smooth] 
        plot (\x, {\x * sqrt(max(0, 1 - 0.5*\x*\x))});
    \draw[very thick, red!80!black, domain=0:1.414, samples=15, smooth] 
        plot (\x, {-\x * sqrt(max(0, 1 - 0.5*\x*\x))});
    % Lóbulo izquierdo (x in [-1.414, 0])
    \draw[very thick, red!80!black, domain=-1.414:0, samples=15, smooth] 
        plot (\x, {\x * sqrt(max(0, 1 - 0.5*\x*\x))});
    \draw[very thick, red!80!black, domain=-1.414:0, samples=15, smooth] 
        plot (\x, {-\x * sqrt(max(0, 1 - 0.5*\x*\x))});

    % Flechas direccionales sobre la separatriz homoclínica
    % Lóbulo derecho: sentido horario
    \draw[-{Latex[length=2.5mm]}, red!80!black, very thick] (0.85, 0.70) -- (1.10, 0.69);
    \draw[-{Latex[length=2.5mm]}, red!80!black, very thick] (1.10, -0.69) -- (0.85, -0.70);
    % Lóbulo izquierdo: sentido horario
    \draw[-{Latex[length=2.5mm]}, red!80!black, very thick] (-1.10, 0.69) -- (-0.85, 0.70);
    \draw[-{Latex[length=2.5mm]}, red!80!black, very thick] (-0.85, -0.70) -- (-1.10, -0.69);

    % Puntos de equilibrio
    % Punto Silla en (0, 0)
    \draw[black, thick, fill=white] (0, 0) circle (3pt);
    \draw[red!80!black, thick] (-0.08, -0.08) -- (0.08, 0.08);
    \draw[red!80!black, thick] (-0.08, 0.08) -- (0.08, -0.08);
    \node[above, font=\scriptsize\bfseries, red!80!black] at (0, 0.12) {Silla $(0, 0)$};

    % Centros en (+-1, 0)
    \filldraw[unmsmblue] (1, 0) circle (2.2pt) node[above right, font=\scriptsize\bfseries, unmsmblue] {Centro $(1, 0)$};
    \filldraw[unmsmblue] (-1, 0) circle (2.2pt) node[above left, font=\scriptsize\bfseries, unmsmblue] {Centro $(-1, 0)$};

    % Etiquetas descriptivas
    \node[above, font=\footnotesize\bfseries, red!80!black] at (1.2, 0.76) {Separatriz Homoclínica ($H = 0$)};
    \node[below, font=\tiny, cyan!75!blue] at (1.0, -0.18) {$H = -\frac{1}{4}$};
    \node[below, font=\tiny, cyan!75!blue] at (-1.0, -0.18) {$H = -\frac{1}{4}$};
\end{tikzpicture}
\end{center}

\end{solucion}

\vspace{3mm}

\begin{problema}[Problema 2.5]
Demuestre que el siguiente sistema dinámico posee al menos una órbita periódica cerrada (ciclo límite) no trivial:
\[
\dot{x} = x - y - x(x^2 + y^2) + \frac{x y}{\sqrt{x^2+y^2}}, \qquad \dot{y} = x + y - y(x^2 + y^2) - \frac{x^2}{\sqrt{x^2+y^2}}
\]
\end{problema}

\begin{solucion}

\textbf{\large 1. Transformación a coordenadas polares}
\[
x = r\cos\theta, \qquad y = r\sin\theta, \qquad r^2 = x^2 + y^2
\]
\[
r\dot{r} = x\dot{x} + y\dot{y}, \qquad r^2\dot{\theta} = x\dot{y} - y\dot{x}
\]

\textbf{Deducción de la ecuación radial:}
\[
x\dot{x} + y\dot{y} = x\left[x - y - x r^2 + \frac{xy}{r}\right] + y\left[x + y - y r^2 - \frac{x^2}{r}\right]
\]
\[
= (x^2 + y^2) - (xy - xy) - r^2(x^2 + y^2) + \frac{x^2y - x^2y}{r} = r^2 - r^4
\]
\[
r\dot{r} = r^2(1 - r^2) \implies \boxed{\dot{r} = r(1 - r^2)}
\]

\textbf{Deducción de la ecuación angular:}
\[
x\dot{y} - y\dot{x} = x\left[x + y - y r^2 - \frac{x^2}{r}\right] - y\left[x - y - x r^2 + \frac{xy}{r}\right]
\]
\[
= (x^2 + y^2) + (xy - xy) - r^2(xy - xy) - \frac{x(x^2 + y^2)}{r} = r^2 - \frac{x r^2}{r} = r^2 - xr
\]
\[
r^2\dot{\theta} = r^2 - r^2\cos\theta = r^2(1 - \cos\theta) \implies \boxed{\dot{\theta} = 1 - \cos\theta \ge 0}
\]
\[
\dot{\theta} > 0 \quad \forall \theta \neq 2k\pi \implies \text{Rotación global en sentido antihorario}
\]

\vspace{3.5mm}

\textbf{\large 2. Puntos de equilibrio del sistema}
\[
\dot{r} = 0 \iff r = 0 \quad \lor \quad r = 1
\]
\[
\dot{\theta} = 0 \iff 1 - \cos\theta = 0 \iff \theta = 2k\pi
\]
\[
\text{En } r = 0: \quad (x, y) = (0, 0) \implies \boxed{P_0 = (0, 0)}
\]

\textbf{Linealización en el origen $(0, 0)$:}
\[
J(0, 0) = \begin{pmatrix} 1 & -1 \\ 1 & 1 \end{pmatrix}
\]
\[
\det(J - \lambda I) = (1 - \lambda)^2 + 1 = 0 \implies \lambda_{1,2} = 1 \pm i
\]
\[
\boxed{\mathrm{Re}(\lambda) = 1 > 0 \implies \text{Foco Inestable (Fuente Repulsora)}}
\]

\vspace{3.5mm}

\textbf{\large 3. Construcción del anillo invariante (región trampa $K$)}
\[
K = \left\{(x, y) \in \mathbb{R}^2 : r_1 \le r \le r_2\right\} = \left\{(x, y) \in \mathbb{R}^2 : \frac{1}{2} \le r \le 2\right\}
\]

\textbf{Comportamiento en la frontera interior $\partial K_{\text{int}} \, (r = 1/2)$:}
\[
\left.\dot{r}\right|_{r=1/2} = \frac{1}{2}\left(1 - \frac{1}{4}\right) = \frac{3}{8} > 0 \implies \boxed{\text{El flujo apunta hacia el exterior de la bola } B_{1/2}(0) \text{ (entra a } K)}
\]

\textbf{Comportamiento en la frontera exterior $\partial K_{\text{ext}} \, (r = 2)$:}
\[
\left.\dot{r}\right|_{r=2} = 2(1 - 4) = -6 < 0 \implies \boxed{\text{El flujo apunta hacia el interior de la bola } B_2(0) \text{ (entra a } K)}
\]

\textbf{Propiedades topológicas de la región trampa $K$:}
\[
K \text{ es cerrado y acotado} \implies \boxed{K \text{ es compacto en } \mathbb{R}^2}
\]
\[
\forall \mathbf{x}_0 \in K, \quad \phi_t(\mathbf{x}_0) \in K \quad \forall t \ge 0 \implies \boxed{K \text{ es positivamente invariante}}
\]
\[
(0, 0) \notin K \implies \boxed{K \text{ no contiene puntos de equilibrio}}
\]

\vspace{3.5mm}

\textbf{\large 4. Aplicación del Teorema de Poincaré-Bendixson}
\[
\boxed{\text{Existe al menos una órbita periódica cerrada (ciclo límite) } \gamma \subset K}
\]

\textbf{Ecuación y estabilidad analítica del ciclo límite:}
\[
\dot{r} = 0 \iff r(1 - r^2) = 0 \implies \boxed{r^* = 1 \iff x^2 + y^2 = 1}
\]
\[
\left.\frac{d\dot{r}}{dr}\right|_{r=1} = \left.(1 - 3r^2)\right|_{r=1} = 1 - 3 = -2 < 0
\]
\[
\boxed{r^* = 1 \implies \text{Ciclo Límite Asintóticamente Estable (Atractor)}}
\]

\vspace{3.5mm}

\textbf{\large 5. Retrato de fase y dinámica orbital en el plano}

\begin{center}
\begin{tikzpicture}[scale=1.25]
    % Región trampa K sombreada (anillo entre r=0.5 y r=2.0)
    \fill[solgreen!35, opacity=0.7] (0, 0) circle (2.0cm);
    \fill[white] (0, 0) circle (0.5cm);

    % Ejes coordenados
    \draw[->, thick, gray!70] (-2.5, 0) -- (2.6, 0) node[right, black, font=\footnotesize] {$x$};
    \draw[->, thick, gray!70] (0, -2.5) -- (0, 2.6) node[above, black, font=\footnotesize] {$y$};

    % Graduaciones en los ejes
    \draw (1, 0.05) -- (1, -0.05) node[below, font=\scriptsize] {$1$};
    \draw (-1, 0.05) -- (-1, -0.05) node[below, font=\scriptsize] {$-1$};
    \draw (2, 0.05) -- (2, -0.05) node[below, font=\scriptsize] {$2$};
    \draw (-2, 0.05) -- (-2, -0.05) node[below, font=\scriptsize] {$-2$};
    \draw (0.05, 1) -- (-0.05, 1) node[left, font=\scriptsize] {$1$};
    \draw (0.05, -1) -- (-0.05, -1) node[left, font=\scriptsize] {$-1$};

    % Fronteras de la región trampa K
    \draw[dashed, thick, deepteal] (0, 0) circle (0.5cm);
    \node[deepteal, font=\tiny\bfseries, below left] at (-0.35, -0.35) {$r = 1/2$};
    \draw[dashed, thick, deepteal] (0, 0) circle (2.0cm);
    \node[deepteal, font=\tiny\bfseries, above right] at (1.4, 1.4) {$r = 2$};

    % Flechas radiales en las fronteras de K indicando flujo entrante
    % Frontera interior r=0.5 (entra a K saliendo del centro)
    \draw[-{Latex[length=2.0mm]}, deepteal, very thick] (0.5, 0) -- (0.78, 0);
    \draw[-{Latex[length=2.0mm]}, deepteal, very thick] (-0.5, 0) -- (-0.78, 0);
    \draw[-{Latex[length=2.0mm]}, deepteal, very thick] (0, 0.5) -- (0, 0.78);
    \draw[-{Latex[length=2.0mm]}, deepteal, very thick] (0, -0.5) -- (0, -0.78);

    % Frontera exterior r=2.0 (entra a K convergiendo hacia el centro)
    \draw[-{Latex[length=2.0mm]}, deepteal, very thick] (2.0, 0) -- (1.70, 0);
    \draw[-{Latex[length=2.0mm]}, deepteal, very thick] (-2.0, 0) -- (-1.70, 0);
    \draw[-{Latex[length=2.0mm]}, deepteal, very thick] (0, 2.0) -- (0, 1.70);
    \draw[-{Latex[length=2.0mm]}, deepteal, very thick] (0, -2.0) -- (0, -1.70);

    % Ciclo límite r* = 1 (línea continua gruesa verde)
    \draw[very thick, solborder] (0, 0) circle (1.0cm);
    \draw[-{Latex[length=2.5mm]}, solborder, very thick] (0, 1.0) -- (-0.2, 1.0);
    \draw[-{Latex[length=2.5mm]}, solborder, very thick] (0, -1.0) -- (0.2, -1.0);
    \node[above, font=\footnotesize\bfseries, solborder] at (0, 1.05) {Ciclo Límite $r^* = 1$};

    % Trayectoria interior que espirala hacia el ciclo límite
    \draw[thick, unmsmblue, -{Latex[length=2.2mm]}] 
        (0.2, 0.1) to[out=70, in=200] (0.55, 0.6) to[out=20, in=90] (0.9, 0.1) 
        to[out=-90, in=0] (0.1, -0.92) to[out=180, in=-90] (-0.95, 0) to[out=90, in=170] (-0.1, 0.98);

    % Trayectoria exterior que espirala hacia el ciclo límite
    \draw[thick, darkpurple, -{Latex[length=2.2mm]}] 
        (1.9, 1.2) to[out=140, in=45] (0.3, 1.7) to[out=-135, in=90] (-1.4, 0.5) 
        to[out=-90, in=180] (-0.3, -1.35) to[out=0, in=-90] (1.2, -0.3) to[out=90, in=-20] (0.9, 0.6);

    % Origen (Foco inestable)
    \draw[black, thick, fill=white] (0, 0) circle (2.8pt);
    \draw[red!80!black, thick] (-0.07, -0.07) -- (0.07, 0.07);
    \draw[red!80!black, thick] (-0.07, 0.07) -- (0.07, -0.07);
    \node[below right, font=\scriptsize\bfseries, red!80!black] at (0.05, -0.05) {$(0,0)$ Fuente};

    % Leyenda informativa
    \node[draw=solborder, fill=white, rounded corners=1mm, font=\tiny\bfseries, align=left] at (1.75, -1.8) {
        \textcolor{solborder}{\textbf{Región trampa }} $K: \frac{1}{2} \le r \le 2$\\
        \textcolor{unmsmblue}{$\bullet$ Flujo interior: } $\dot{r} > 0$\\
        \textcolor{darkpurple}{$\bullet$ Flujo exterior: } $\dot{r} < 0$\\
        \textcolor{solborder}{$\bullet$ Atractor: } $r^* = 1$
    };
\end{tikzpicture}
\end{center}

\end{solucion}

\begin{conclusion}
La existencia del ciclo límite queda rigurosamente establecida por el Teorema de Poincaré-Bendixson: el anillo cerrado $K$ confina todas las trayectorias que entran en él, y al no existir puntos fijos en su interior, el flujo queda atrapado asintóticamente en una oscilación auto-sostenida estable (ciclo límite $r^* = 1$).
\end{conclusion}

\vspace{5mm}
% =========================================================================
% TEMA 3: TEORÍA DE BIFURCACIONES
% =========================================================================
\section*{\color{unmsmblue} Tema 3: Teoría de Bifurcaciones}

\begin{problema}[Problema 3.1]
Estudie los puntos fijos, su estabilidad y dibuje el diagrama de bifurcación para el sistema:
\[
\dot{x} = \mu - x^2, \quad \mu \in \mathbb{R}
\]
\end{problema}

\begin{solucion}

\textbf{\large 1. Determinar los puntos de equilibrio en función del parámetro $\mu$}
\[
f(x; \mu) = \mu - x^2 = 0 \iff x^2 = \mu
\]

\textbf{Caso $\mu < 0$:}
\[
x^2 = \mu < 0 \implies \boxed{\nexists \, x^* \in \mathbb{R} \quad (\text{No existen puntos de equilibrio reales})}
\]
\[
\dot{x} = \mu - x^2 \le \mu < 0 \quad \forall x \in \mathbb{R} \implies \text{El flujo decae monótonamente hacia } -\infty
\]

\textbf{Caso $\mu = 0$ (Punto crítico de bifurcación):}
\[
x^2 = 0 \implies \boxed{\mu_c = 0 \implies x^* = 0 \quad (\text{Equilibrio no hiperbólico semistable})}
\]

\textbf{Caso $\mu > 0$:}
\[
x^2 = \mu > 0 \implies \boxed{x_1^*(\mu) = +\sqrt{\mu}, \qquad x_2^*(\mu) = -\sqrt{\mu}}
\]

\vspace{3.5mm}

\textbf{\large 2. Análisis de estabilidad lineal}
\[
f'(x; \mu) = \frac{\partial}{\partial x}\left(\mu - x^2\right) = -2x
\]

\textbf{Evaluación en la rama superior $x_1^*(\mu) = +\sqrt{\mu}$:}
\[
f'(+\sqrt{\mu}) = -2\sqrt{\mu} < 0 \quad (\forall \mu > 0)
\]
\[
\boxed{x_1^*(\mu) = +\sqrt{\mu} \implies \text{Punto Fijo Asintóticamente Estable (Nodo Atractor)}}
\]

\textbf{Evaluación en la rama inferior $x_2^*(\mu) = -\sqrt{\mu}$:}
\[
f'(-\sqrt{\mu}) = -2(-\sqrt{\mu}) = +2\sqrt{\mu} > 0 \quad (\forall \mu > 0)
\]
\[
\boxed{x_2^*(\mu) = -\sqrt{\mu} \implies \text{Punto Fijo Inestable (Nodo Repulsor)}}
\]

\textbf{Evaluación en el punto crítico $(\mu_c, x^*) = (0, 0)$:}
\[
f'(0; 0) = 0 \implies \boxed{\lambda = 0 \quad (\text{Pérdida de hiperbolicidad}) }
\]

\vspace{3.5mm}

\textbf{\large 3. Verificación de las condiciones de no degeneración de Sotomayor}
\[
f(0; 0) = 0 - 0^2 = 0
\]
\[
f_x(0; 0) = -2(0) = 0
\]
\[
f_\mu(0; 0) = \left.\frac{\partial}{\partial \mu}(\mu - x^2)\right|_{(0,0)} = 1 \neq 0 \quad (\text{Transversalidad del parámetro})
\]
\[
f_{xx}(0; 0) = \left.\frac{\partial^2}{\partial x^2}(\mu - x^2)\right|_{(0,0)} = -2 \neq 0 \quad (\text{No degeneración cuadrática})
\]
\[
\boxed{\text{Se satisfacen rigurosamente las hipótesis de la Bifurcación Silla-Nodo genérica en } (\mu_c, x^*) = (0, 0)}
\]

\vspace{3.5mm}

\textbf{\large 4. Diagrama de bifurcación y flujo en el espacio de fases}

\begin{center}
\begin{tikzpicture}[scale=1.35]
    % Zona sin equilibrios sombreada (mu < 0)
    \fill[warnbg!55, opacity=0.5] (-1.6, -2.1) rectangle (0, 2.1);
    \node[noteborder, font=\tiny\bfseries, align=center] at (-0.8, 1.45) {Sin equilibrios\\$(\mu < 0)$\\$\dot{x} < 0$};

    % Ejes coordenados del diagrama de bifurcación
    \draw[->, thick, gray!70] (-1.7, 0) -- (3.2, 0) node[right, black, font=\footnotesize] {Parámetro $\mu$};
    \draw[->, thick, gray!70] (0, -2.2) -- (0, 2.3) node[above, black, font=\footnotesize] {Estado $x^*$};

    % Graduaciones en los ejes
    \draw (1, 0.05) -- (1, -0.05) node[below, font=\scriptsize] {$1$};
    \draw (2, 0.05) -- (2, -0.05) node[below, font=\scriptsize] {$2$};
    \draw (-1, 0.05) -- (-1, -0.05) node[below, font=\scriptsize] {$-1$};
    \draw (0.05, 1) -- (-0.05, 1) node[left, font=\scriptsize] {$1$};
    \draw (0.05, -1) -- (-0.05, -1) node[left, font=\scriptsize] {$-1$};

    % Rama estable: x* = +sqrt(mu) (línea continua verde)
    \draw[very thick, solborder, domain=0:2.6, samples=16, smooth] plot (\x, {sqrt(\x)});
    \node[right, font=\scriptsize\bfseries, solborder] at (2.6, 1.62) {Rama Estable $x_1^* = +\sqrt{\mu}$};

    % Rama inestable: x* = -sqrt(mu) (línea discontinua roja)
    \draw[very thick, dashed, red!80!black, domain=0:2.6, samples=16, smooth] plot (\x, {-sqrt(\x)});
    \node[right, font=\scriptsize\bfseries, red!80!black] at (2.6, -1.62) {Rama Inestable $x_2^* = -\sqrt{\mu}$};

    % Vectores de flujo vertical para mu < 0 (flujo continuo hacia abajo)
    \draw[-{Latex[length=2.0mm]}, thick, gray!80!black] (-0.8, 1.9) -- (-0.8, 1.35);
    \draw[-{Latex[length=2.0mm]}, thick, gray!80!black] (-0.8, 0.3) -- (-0.8, -0.3);
    \draw[-{Latex[length=2.0mm]}, thick, gray!80!black] (-0.8, -1.35) -- (-0.8, -1.9);

    % Vectores de flujo vertical para mu = 1.5 > 0
    % Superior: por encima de x1* (apunta hacia x1*)
    \draw[-{Latex[length=2.0mm]}, thick, deepteal] (1.5, 2.05) -- (1.5, 1.45);
    % Intermedio: entre x2* y x1* (apunta hacia x1*)
    \draw[-{Latex[length=2.0mm]}, thick, deepteal] (1.5, -0.7) -- (1.5, 0.7);
    % Inferior: por debajo de x2* (repelido hacia -infinito)
    \draw[-{Latex[length=2.0mm]}, thick, deepteal] (1.5, -1.45) -- (1.5, -2.05);

    % Punto de bifurcación crítico en el origen
    \filldraw[unmsmblue] (0, 0) circle (3pt);
    \node[above left, font=\scriptsize\bfseries, unmsmblue] at (-0.05, 0.08) {Bifurcación Silla-Nodo $(\mu_c = 0)$};
\end{tikzpicture}
\end{center}

\end{solucion}

\begin{conclusion}
La bifurcación silla-nodo representa el mecanismo fundamental de creación y aniquilación de estados estacionarios. Al variar $\mu$ desde valores negativos hacia positivos, dos equilibrios de estabilidad opuesta colisionan en $\mu_c = 0$ y se bifurcan como una parábola orientada hacia la derecha en el espacio $(\mu, x^*)$.
\end{conclusion}

\vspace{3mm}

\begin{problema}[Problema 3.2]
Analice la bifurcación del sistema $\dot{x} = \mu x - x^2$ y determine el intercambio de estabilidad.
\end{problema}

\begin{solucion}

\textbf{\large 1. Determinar los puntos de equilibrio del sistema}
\[
f(x; \mu) = \mu x - x^2 = x(\mu - x) = 0
\]
\[
\boxed{x_1^*(\mu) = 0, \qquad x_2^*(\mu) = \mu \quad (\forall \mu \in \mathbb{R})}
\]

\textbf{Punto de colisión e intersección transversal:}
\[
x_1^*(\mu) = x_2^*(\mu) \iff \mu = 0 \implies \boxed{\mu_c = 0 \implies x^* = 0}
\]

\vspace{3.5mm}

\textbf{\large 2. Análisis de estabilidad lineal}
\[
f'(x; \mu) = \frac{\partial}{\partial x}\left(\mu x - x^2\right) = \mu - 2x
\]

\textbf{Evaluación en la rama horizontal trivial $x_1^*(\mu) = 0$:}
\[
f'(0; \mu) = \mu - 2(0) = \mu
\]
\[
\mu < 0 \implies f'(0) < 0 \implies \boxed{x_1^* = 0 \implies \text{Asintóticamente Estable (Atractor)}}
\]
\[
\mu > 0 \implies f'(0) > 0 \implies \boxed{x_1^* = 0 \implies \text{Inestable (Repulsor)}}
\]

\textbf{Evaluación en la rama diagonal móvil $x_2^*(\mu) = \mu$:}
\[
f'(\mu; \mu) = \mu - 2\mu = -\mu
\]
\[
\mu < 0 \implies f'(\mu) = -\mu > 0 \implies \boxed{x_2^* = \mu \implies \text{Inestable (Repulsor)}}
\]
\[
\mu > 0 \implies f'(\mu) = -\mu < 0 \implies \boxed{x_2^* = \mu \implies \text{Asintóticamente Estable (Atractor)}}
\]

\textbf{Evaluación en el punto crítico $(\mu_c, x^*) = (0, 0)$:}
\[
f'(0; 0) = 0 \implies \boxed{\lambda = 0 \quad (\text{Pérdida de hiperbolicidad y cruce de autovalores})}
\]

\vspace{3.5mm}

\textbf{\large 3. Verificación de las condiciones de Sotomayor para bifurcación transcrítica}
\[
f(0; 0) = 0(0 - 0) = 0
\]
\[
f_x(0; 0) = \left.(\mu - 2x)\right|_{(0,0)} = 0
\]
\[
f_\mu(0; 0) = \left.x\right|_{(0,0)} = 0 \quad (\text{Garantiza la persistencia del equilibrio trivial } x^* = 0)
\]
\[
f_{x\mu}(0; 0) = \left.\frac{\partial^2 f}{\partial x \partial \mu}\right|_{(0,0)} = 1 \neq 0 \quad (\text{Condición de transversalidad en el cambio de signo de } \lambda)
\]
\[
f_{xx}(0; 0) = \left.\frac{\partial^2 f}{\partial x^2}\right|_{(0,0)} = -2 \neq 0 \quad (\text{Curvatura cuadrática no nula})
\]
\[
\boxed{\text{Se satisfacen rigurosamente las hipótesis de la Bifurcación Transcrítica genérica en } (\mu_c, x^*) = (0, 0)}
\]

\vspace{3.5mm}

\textbf{\large 4. Diagrama de bifurcación y flujo en el espacio de fases}

\begin{center}
\begin{tikzpicture}[scale=1.4]
    % Fondo sombreado sutil para diferenciar regiones de estabilidad
    \fill[solgreen!20, opacity=0.5] (-2.2, 0) -- (0, 0) -- (-2.2, -2.2) -- cycle;
    \fill[solgreen!20, opacity=0.5] (0, 0) -- (2.2, 0) -- (2.2, 2.2) -- cycle;

    % Ejes coordenados del diagrama de bifurcación
    \draw[->, thick, gray!70] (-2.3, 0) -- (2.5, 0) node[right, black, font=\footnotesize] {Parámetro $\mu$};
    \draw[->, thick, gray!70] (0, -2.1) -- (0, 2.3) node[above, black, font=\footnotesize] {Estado $x^*$};

    % Graduaciones en los ejes
    \draw (1, 0.05) -- (1, -0.05) node[below, font=\scriptsize] {$1$};
    \draw (2, 0.05) -- (2, -0.05) node[below, font=\scriptsize] {$2$};
    \draw (-1, 0.05) -- (-1, -0.05) node[below, font=\scriptsize] {$-1$};
    \draw (-2, 0.05) -- (-2, -0.05) node[below, font=\scriptsize] {$-2$};
    \draw (0.05, 1) -- (-0.05, 1) node[left, font=\scriptsize] {$1$};
    \draw (0.05, -1) -- (-0.05, -1) node[left, font=\scriptsize] {$-1$};

    % Rama x* = 0: Estable para mu < 0, Inestable para mu > 0
    \draw[line width=1.5pt, solborder] (-2.2, 0) -- (0, 0);
    \draw[line width=1.5pt, dashed, red!80!black] (0, 0) -- (2.2, 0);

    % Rama x* = mu: Inestable para mu < 0, Estable para mu > 0
    \draw[line width=1.5pt, dashed, red!80!black] (-1.9, -1.9) -- (0, 0);
    \draw[line width=1.5pt, solborder] (0, 0) -- (1.9, 1.9);

    % Etiquetas de las ramas
    \node[left, font=\scriptsize\bfseries, solborder] at (-1.2, 0.22) {Rama Estable ($x_1^* = 0$)};
    \node[right, font=\scriptsize\bfseries, red!80!black] at (1.1, 0.22) {Rama Inestable ($x_1^* = 0$)};
    \node[above left, font=\scriptsize\bfseries, solborder] at (1.65, 1.7) {Rama Estable ($x_2^* = \mu$)};
    \node[below right, font=\scriptsize\bfseries, red!80!black] at (-1.65, -1.7) {Rama Inestable ($x_2^* = \mu$)};

    % Vectores de flujo vertical en el espacio de fases
    % Para mu = -1.2 < 0:
    \draw[-{Latex[length=1.8mm]}, thick, deepteal] (-1.2, 1.5) -- (-1.2, 0.4);
    \draw[-{Latex[length=1.8mm]}, thick, deepteal] (-1.2, -0.8) -- (-1.2, -0.2);
    \draw[-{Latex[length=1.8mm]}, thick, deepteal] (-1.2, -1.4) -- (-1.2, -1.9);

    % Para mu = +1.2 > 0:
    \draw[-{Latex[length=1.8mm]}, thick, deepteal] (1.2, 1.9) -- (1.2, 1.4);
    \draw[-{Latex[length=1.8mm]}, thick, deepteal] (1.2, 0.2) -- (1.2, 0.8);
    \draw[-{Latex[length=1.8mm]}, thick, deepteal] (1.2, -0.4) -- (1.2, -1.5);

    % Punto crítico de colisión e intercambio
    \filldraw[unmsmblue] (0, 0) circle (3pt);
    \node[below right, font=\scriptsize\bfseries, unmsmblue] at (0.05, -0.08) {Intercambio $(\mu_c = 0)$};
\end{tikzpicture}
\end{center}

\end{solucion}

\begin{conclusion}
La bifurcación transcrítica se distingue porque ambos equilibrios persisten a ambos lados del valor crítico $\mu_c = 0$. Al atravesar el origen, las trayectorias de equilibrio no se destruyen sino que colisionan e intercambian su estabilidad: el equilibrio que atraía el flujo pasa a repelerlo, transfiriendo la estabilidad a la rama que antes era inestable.
\end{conclusion}

\vspace{3mm}

\begin{problema}[Problema 3.3]
Compare la naturaleza de las bifurcaciones:
\[
\text{(A) } \dot{x} = \mu x - x^3 \qquad \text{y} \qquad \text{(B) } \dot{x} = \mu x + x^3 - x^5
\]
\end{problema}

\begin{solucion}

\textbf{\large 1. Análisis de la bifurcación pitchfork supercrítica: $\dot{x} = \mu x - x^3$}
\[
f_A(x; \mu) = \mu x - x^3 = x(\mu - x^2) = 0
\]
\[
\boxed{f_A(-x; \mu) = -f_A(x; \mu) \implies \text{Invarianza bajo la simetría de inversión } x \to -x}
\]

\textbf{Puntos de equilibrio:}
\[
\mu \le 0 \implies \boxed{x^* = 0 \quad (\text{Único punto de equilibrio real})}
\]
\[
\mu > 0 \implies \boxed{x_0^* = 0, \qquad x_{\pm}^* = \pm\sqrt{\mu}}
\]

\textbf{Estabilidad lineal:}
\[
f_A'(x; \mu) = \mu - 3x^2
\]
\[
f_A'(0; \mu) = \mu \implies 
\begin{cases}
\mu < 0 \implies \boxed{x_0^* = 0 \text{ es Asintóticamente Estable}} \\
\mu > 0 \implies \boxed{x_0^* = 0 \text{ es Inestable}}
\end{cases}
\]
\[
f_A'(\pm\sqrt{\mu}; \mu) = \mu - 3(\pm\sqrt{\mu})^2 = \mu - 3\mu = -2\mu < 0 \quad (\forall \mu > 0)
\]
\[
\boxed{x_{\pm}^* = \pm\sqrt{\mu} \implies \text{Ramas Simétricas Asintóticamente Estables (Atractores)}}
\]

\vspace{3.5mm}

\textbf{\large 2. Análisis de la bifurcación pitchfork subcrítica estabilizada: $\dot{x} = \mu x + x^3 - x^5$}
\[
f_B(x; \mu) = x(\mu + x^2 - x^4) = 0
\]

\textbf{Equilibrio trivial:}
\[
x_0^* = 0 \quad \forall \mu \in \mathbb{R}, \qquad f_B'(0; \mu) = \mu
\]
\[
\begin{cases}
\mu < 0 \implies \boxed{x_0^* = 0 \text{ es Asintóticamente Estable}} \\
\mu > 0 \implies \boxed{x_0^* = 0 \text{ es Inestable}}
\end{cases}
\]

\textbf{Ramas no triviales ($u = x^2 > 0$):}
\[
u^2 - u - \mu = 0 \implies u = x^2 = \frac{1 \pm \sqrt{1 + 4\mu}}{2}
\]
\[
1 + 4\mu \ge 0 \iff \boxed{\mu \ge -\frac{1}{4} = \mu_s \quad (\text{Punto de bifurcación silla-nodo de ramas laterales})}
\]

\textbf{Rango de biestabilidad e histéresis $\mu \in (-1/4, \, 0)$:}
\[
\boxed{x_{\text{int}}^* = \pm\sqrt{\frac{1 - \sqrt{1 + 4\mu}}{2}} \quad (\text{Ramas intermedias inestables})}
\]
\[
\boxed{x_{\text{ext}}^* = \pm\sqrt{\frac{1 + \sqrt{1 + 4\mu}}{2}} \quad (\text{Ramas exteriores de gran amplitud estables})}
\]

\textbf{Rango post-crítico $\mu > 0$:}
\[
\boxed{x_{\text{ext}}^* = \pm\sqrt{\frac{1 + \sqrt{1 + 4\mu}}{2}} \quad (\text{Únicos atractores estables existentes})}
\]

\vspace{3.5mm}

\textbf{\large 3. Fenómenos dinámicos característicos: Catástrofe, salto de estado e histéresis}
\[
\mu \to 0^+ \implies \text{El estado en } x = 0 \text{ colapsa y salta discontinuamente a } \boxed{x^* = \pm 1}
\]
\[
\mu \to \mu_s^+ = -1/4 \implies \text{Las ramas exteriores se aniquilan con las inestables y colapsan a } \boxed{x^* = 0}
\]
\[
\boxed{\text{Ciclo de Histéresis: } \Delta\mu = 0 - \left(-\frac{1}{4}\right) = \frac{1}{4}}
\]

\vspace{3.5mm}

\textbf{\large 4. Comparación de los diagramas de bifurcación}

\begin{center}
\begin{tikzpicture}[scale=1.15]
    % ================= PANEL A: SUPERCRÍTICA =================
    \begin{scope}[xshift=-3.6cm]
        \node[above, font=\footnotesize\bfseries, unmsmblue] at (0, 2.2) {Caso (A): Supercrítica ($\dot{x} = \mu x - x^3$)};

        % Ejes
        \draw[->, thick, gray!70] (-1.8, 0) -- (2.3, 0) node[right, black, font=\scriptsize] {$\mu$};
        \draw[->, thick, gray!70] (0, -2.0) -- (0, 2.1) node[above, black, font=\scriptsize] {$x^*$};

        % Graduaciones
        \draw (1, 0.05) -- (1, -0.05) node[below, font=\tiny] {$1$};
        \draw (-1, 0.05) -- (-1, -0.05) node[below, font=\tiny] {$-1$};
        \draw (0.05, 1) -- (-0.05, 1) node[left, font=\tiny] {$1$};
        \draw (0.05, -1) -- (-0.05, -1) node[left, font=\tiny] {$-1$};

        % Rama trivial x* = 0
        \draw[line width=1.5pt, solborder] (-1.7, 0) -- (0, 0);
        \draw[line width=1.5pt, dashed, red!80!black] (0, 0) -- (2.1, 0);

        % Ramas x* = +-sqrt(mu)
        \draw[line width=1.5pt, solborder, domain=0:2.1, samples=15, smooth] plot (\x, {sqrt(\x)});
        \draw[line width=1.5pt, solborder, domain=0:2.1, samples=15, smooth] plot (\x, {-sqrt(\x)});

        % Punto de bifurcación
        \filldraw[unmsmblue] (0, 0) circle (2.8pt);
        \node[below left, font=\tiny\bfseries, unmsmblue] at (0, -0.05) {$\mu_c = 0$};

        % Etiquetas
        \node[right, font=\tiny\bfseries, solborder] at (2.1, 1.45) {$+\sqrt{\mu}$};
        \node[right, font=\tiny\bfseries, solborder] at (2.1, -1.45) {$-\sqrt{\mu}$};
        \node[above, font=\tiny\bfseries, solborder] at (-0.9, 0.05) {Estable};
        \node[above, font=\tiny\bfseries, red!80!black] at (1.0, 0.05) {Inestable};
        \node[draw=solborder, fill=white, rounded corners=0.8mm, font=\tiny\bfseries, solborder] at (0, -1.7) {Transición continua (segundo orden)};
    \end{scope}

    % ================= PANEL B: SUBCRÍTICA =================
    \begin{scope}[xshift=3.6cm]
        \node[above, font=\footnotesize\bfseries, darkpurple] at (0, 2.2) {Caso (B): Subcrítica ($\dot{x} = \mu x + x^3 - x^5$)};

        % Zona de biestabilidad sombreada
        \fill[solgreen!25, opacity=0.6] (-0.6, -1.9) rectangle (0, 1.9);

        % Ejes
        \draw[->, thick, gray!70] (-1.8, 0) -- (2.3, 0) node[right, black, font=\scriptsize] {$\mu$};
        \draw[->, thick, gray!70] (0, -2.0) -- (0, 2.1) node[above, black, font=\scriptsize] {$x^*$};

        % Graduaciones
        \draw (1, 0.05) -- (1, -0.05) node[below, font=\tiny] {$1$};
        \draw (-1, 0.05) -- (-1, -0.05) node[below, font=\tiny] {$-1$};
        \draw (-0.6, 0.05) -- (-0.6, -0.05) node[below, font=\tiny] {$\mu_s$};

        % Rama trivial x* = 0
        \draw[line width=1.5pt, solborder] (-1.7, 0) -- (0, 0);
        \draw[line width=1.5pt, dashed, red!80!black] (0, 0) -- (2.1, 0);

        % Ramas intermedias inestables (curvan hacia atras)
        \draw[line width=1.4pt, dashed, red!80!black] (0, 0) to[out=130, in=90] (-0.6, 0.7);
        \draw[line width=1.4pt, dashed, red!80!black] (0, 0) to[out=-130, in=-90] (-0.6, -0.7);

        % Ramas exteriores estables (desde mu_s hacia adelante)
        \draw[line width=1.5pt, solborder] (-0.6, 0.7) to[out=-90, in=180] (-0.35, 0.95) to[out=0, in=190] (2.0, 1.6);
        \draw[line width=1.5pt, solborder] (-0.6, -0.7) to[out=90, in=180] (-0.35, -0.95) to[out=0, in=170] (2.0, -1.6);

        % Flechas de salto discontinuo / catástrofe
        \draw[-{Latex[length=2.2mm]}, red!80!black, very thick] (0.05, 0.15) -- (0.05, 1.15);
        \draw[-{Latex[length=2.2mm]}, deepteal, very thick] (-0.55, 0.9) -- (-0.55, 0.15);

        % Puntos de bifurcación
        \filldraw[unmsmblue] (0, 0) circle (2.8pt);
        \filldraw[darkpurple] (-0.6, 0.7) circle (2.2pt);
        \filldraw[darkpurple] (-0.6, -0.7) circle (2.2pt);

        % Etiquetas
        \node[right, font=\tiny\bfseries, solborder] at (2.0, 1.6) {Rama Ext.};
        \node[right, font=\tiny\bfseries, solborder] at (2.0, -1.6) {Rama Ext.};
        \node[draw=darkpurple, fill=white, rounded corners=0.8mm, font=\tiny\bfseries, darkpurple] at (0, -1.7) {Biestabilidad e Histéresis};
    \end{scope}
\end{tikzpicture}
\end{center}

\end{solucion}

\begin{conclusion}
La comparación entre ambas bifurcaciones trinchera evidencia un contraste fundamental en dinámica no lineal:
\begin{enumerate}
    \item La bifurcación supercrítica es segura y reversible (transición continua de segundo orden): al variar $\mu$ a través de cero, la amplitud del nuevo estado crece suavemente como $\sqrt{\mu}$.
    \item La bifurcación subcrítica es explosiva y peligrosa (transición discontinua de primer orden): al cruzar $\mu = 0$, el sistema sufre un salto repentino e irreversible de gran amplitud (catástrofe), requiriendo reducir el parámetro hasta $\mu_s = -1/4$ para retornar al estado base original (histéresis).
\end{enumerate}
\end{conclusion}

\vspace{3mm}

\begin{problema}[Problema 3.4]
Demuestre que el siguiente sistema presenta una bifurcación de Hopf en $\mu = 0$ y determine la amplitud y estabilidad del ciclo límite que nace:
\[
\dot{x} = \mu x - y - x(x^2 + y^2), \qquad \dot{y} = x + \mu y - y(x^2 + y^2)
\]
\end{problema}

\begin{solucion}

\textbf{\large 1. Linealización en el origen y verificación de las condiciones de Hopf}
\[
J(x, y) = \begin{pmatrix} \mu - 3x^2 - y^2 & -1 - 2xy \\ 1 - 2xy & \mu - x^2 - 3y^2 \end{pmatrix}
\]
\[
J(0, 0) = \begin{pmatrix} \mu & -1 \\ 1 & \mu \end{pmatrix}
\]

\textbf{Ecuación característica y autovalores:}
\[
\det(J(0,0) - \lambda I) = (\mu - \lambda)^2 + 1 = 0 \implies \lambda_{1,2}(\mu) = \mu \pm i = \alpha(\mu) \pm i\omega(\mu)
\]
\[
\boxed{\alpha(\mu) = \mathrm{Re}(\lambda) = \mu, \qquad \omega(\mu) = \mathrm{Im}(\lambda) = 1}
\]

\textbf{Verificación de la condición espectral de Hopf en $\mu_c = 0$:}
\[
\alpha(0) = 0, \qquad \omega(0) = \omega_0 = 1 \neq 0 \implies \boxed{\lambda_{1,2}(0) = \pm i \quad (\text{Par de autovalores imaginarios puros})}
\]

\textbf{Verificación de la condición de transversalidad:}
\[
\left.\frac{d\alpha}{d\mu}\right|_{\mu = 0} = \left.\frac{d}{d\mu}(\mu)\right|_{\mu = 0} = 1 \neq 0 \implies \boxed{\text{Cruce transversal a través del eje imaginario}}
\]

\vspace{3.5mm}

\textbf{\large 2. Transformación canónica a coordenadas polares}
\[
x = r\cos\theta, \qquad y = r\sin\theta, \qquad r^2 = x^2 + y^2
\]
\[
r\dot{r} = x\dot{x} + y\dot{y}, \qquad r^2\dot{\theta} = x\dot{y} - y\dot{x}
\]

\textbf{Deducción de la ecuación radial:}
\[
x\dot{x} + y\dot{y} = x\left[\mu x - y - x(x^2 + y^2)\right] + y\left[x + \mu y - y(x^2 + y^2)\right]
\]
\[
= \mu(x^2 + y^2) - (xy - xy) - (x^2 + y^2)^2 = \mu r^2 - r^4
\]
\[
r\dot{r} = r^2(\mu - r^2) \implies \boxed{\dot{r} = \mu r - r^3 = r(\mu - r^2)}
\]

\textbf{Deducción de la ecuación angular:}
\[
x\dot{y} - y\dot{x} = x\left[x + \mu y - y r^2\right] - y\left[\mu x - y - x r^2\right]
\]
\[
= (x^2 + y^2) + \mu(xy - xy) - r^2(xy - xy) = r^2
\]
\[
r^2\dot{\theta} = r^2 \implies \boxed{\dot{\theta} = 1 > 0 \quad (\text{Velocidad angular constante en sentido antihorario})}
\]

\vspace{3.5mm}

\textbf{\large 3. Existencia, radio y estabilidad orbital del ciclo límite}

\textbf{Comportamiento para $\mu \le 0$:}
\[
\dot{r} = r(\mu - r^2) < 0 \quad \forall r > 0 \implies \boxed{r^* = 0 \text{ es el único atractor global (Foco Estable)}}
\]

\textbf{Nacimiento del ciclo límite para $\mu > 0$:}
\[
\dot{r} = 0 \iff r(\mu - r^2) = 0 \implies \boxed{R(\mu) = \sqrt{\mu} \quad (\mu > 0)}
\]

\textbf{Análisis de estabilidad orbital de la órbita periódica:}
\[
\left.\frac{d\dot{r}}{dr}\right|_{r = \sqrt{\mu}} = \left.\frac{d}{dr}(\mu r - r^3)\right|_{r = \sqrt{\mu}} = \mu - 3(\sqrt{\mu})^2 = \mu - 3\mu = -2\mu < 0 \quad (\forall \mu > 0)
\]
\[
\boxed{R(\mu) = \sqrt{\mu} \implies \text{Ciclo Límite Periódico Asintóticamente Estable (Atractor)}}
\]

\textbf{Periodo exacto de la oscilación auto-sostenida:}
\[
T = \int_0^{2\pi} \frac{d\theta}{\dot{\theta}} = \int_0^{2\pi} 1\, d\theta \implies \boxed{T = 2\pi}
\]

\vspace{3.5mm}

\textbf{\large 4. Coeficiente de primer orden de Lyapunov y clasificación de Hopf}
\[
\dot{r} = \mu r - a r^3, \qquad \text{con } a = 1 > 0 \implies \boxed{l_1 = -a = -1 < 0}
\]
\[
\boxed{\text{La bifurcación es rigurosamente Supercrítica (el ciclo límite es atractor y nace para } \mu > 0)}
\]

\vspace{3.5mm}

\textbf{\large 5. Diagrama de bifurcación de Hopf en el espacio de fases}

\begin{center}
\begin{tikzpicture}[scale=1.3]
    % Región previa a la bifurcación (mu < 0) sombreada
    \fill[probhead!45, opacity=0.5] (-2.2, -1.9) rectangle (0, 1.9);
    \node[unmsmblue, font=\tiny\bfseries, align=center] at (-1.1, 1.45) {Foco Estable\\Sin oscilaciones};

    % Ejes coordenados del diagrama de bifurcación
    \draw[->, thick, gray!70] (-2.3, 0) -- (3.2, 0) node[right, black, font=\footnotesize] {Parámetro $\mu$};
    \draw[->, thick, gray!70] (0, -2.0) -- (0, 2.1) node[above, black, font=\footnotesize] {Estado / Amplitud $x, y$};

    % Graduaciones
    \draw (1, 0.05) -- (1, -0.05) node[below, font=\scriptsize] {$1$};
    \draw (2, 0.05) -- (2, -0.05) node[below, font=\scriptsize] {$2$};
    \draw (-1, 0.05) -- (-1, -0.05) node[below, font=\scriptsize] {$-1$};

    % Origen trivial x* = 0
    \draw[line width=1.5pt, solborder] (-2.1, 0) -- (0, 0);
    \draw[line width=1.5pt, dashed, red!80!black] (0, 0) -- (2.8, 0);
    \node[above, font=\tiny\bfseries, solborder] at (-1.1, 0.06) {Foco Estable ($r^* = 0$)};
    \node[below, font=\tiny\bfseries, red!80!black] at (1.4, -0.06) {Foco Inestable};

    % Envolvente parabólica del ciclo límite R = +-sqrt(mu)
    \draw[line width=1.4pt, deepteal, domain=0:2.7, samples=15, smooth] plot (\x, {sqrt(\x)});
    \draw[line width=1.4pt, deepteal, domain=0:2.7, samples=15, smooth] plot (\x, {-sqrt(\x)});
    \node[right, font=\scriptsize\bfseries, deepteal] at (2.7, 1.64) {$R(\mu) = +\sqrt{\mu}$};
    \node[right, font=\scriptsize\bfseries, deepteal] at (2.7, -1.64) {$R(\mu) = -\sqrt{\mu}$};

    % Órbitas periódicas cerradas representadas en cortes transversales mu = cte
    % En mu = 1.0 (R = 1.0)
    \draw[thick, solborder] (1.0, 0) ellipse (0.16cm and 1.0cm);
    \draw[-{Latex[length=1.8mm]}, solborder, thick] (1.16, 0.1) -- (1.16, 0.25);
    \node[above, font=\tiny\bfseries, solborder] at (1.0, 1.05) {$\mu = 1$};

    % En mu = 2.0 (R = 1.414)
    \draw[thick, solborder] (2.0, 0) ellipse (0.20cm and 1.414cm);
    \draw[-{Latex[length=1.8mm]}, solborder, thick] (2.20, 0.1) -- (2.20, 0.25);
    \node[above, font=\tiny\bfseries, solborder] at (2.0, 1.47) {$\mu = 2$};

    % Flechas de convergencia hacia el ciclo límite
    \draw[-{Latex[length=1.8mm]}, darkpurple, thick] (2.0, 0.3) -- (2.0, 0.95);
    \draw[-{Latex[length=1.8mm]}, darkpurple, thick] (2.0, 1.9) -- (2.0, 1.55);

    % Punto de bifurcación de Hopf
    \filldraw[unmsmblue] (0, 0) circle (3pt);
    \node[above left, font=\scriptsize\bfseries, unmsmblue] at (-0.05, 0.12) {Bifurcación de Hopf $(\mu_c = 0)$};
    \node[draw=solborder, fill=white, rounded corners=0.8mm, font=\tiny\bfseries, solborder] at (1.3, -1.6) {Ciclo Límite Estable (Supercrítico)};
\end{tikzpicture}
\end{center}

\end{solucion}

\begin{conclusion}
La bifurcación de Hopf supercrítica modela la transición dinámica donde un estado estacionario pierde estabilidad y transfiere toda la energía a una oscilación auto-sostenida estable (ciclo límite). La amplitud del ciclo crece suavemente con la ley de potencias $R(\mu) = \sqrt{\mu}$, garantizando una transición continua y sin histéresis hacia el régimen oscilatorio.
\end{conclusion}

\vspace{3mm}

\begin{problema}[Problema 3.5]
Considere el sistema en el plano:
\[
\dot{x} = y, \qquad \dot{y} = \mu y + x - x^2 + x y
\]
Explique cómo la variación del parámetro $\mu$ destruye la órbita homoclínica y da lugar a una bifurcación global de ciclo límite (Bifurcación Homoclínica de Silla).
\end{problema}

\begin{solucion}

\textbf{\large 1. Determinar y clasificar los puntos de equilibrio}
\[
\dot{x} = y = 0
\]
\[
\dot{y} = \mu(0) + x - x^2 + x(0) = x(1 - x) = 0 \iff x = 0 \quad \lor \quad x = 1
\]
\[
\boxed{P_1 = (0, 0), \qquad P_2 = (1, 0)}
\]

\textbf{Matriz jacobiana general:}
\[
J(x, y) = \begin{pmatrix} 0 & 1 \\ 1 - 2x + y & \mu + x \end{pmatrix}
\]

\textbf{Análisis en el origen $P_1(0, 0)$:}
\[
J(0, 0) = \begin{pmatrix} 0 & 1 \\ 1 & \mu \end{pmatrix}
\]
\[
\det(J) = -1 < 0 \quad (\forall \mu \in \mathbb{R})
\]
\[
\boxed{\lambda_{u, s} = \frac{\mu \pm \sqrt{\mu^2 + 4}}{2} \implies \text{Punto Silla Hiperbólico para todo } \mu \in \mathbb{R}}
\]
\[
\lambda_u = \frac{\mu + \sqrt{\mu^2 + 4}}{2} > 0 \quad (\text{Rama inestable}), \qquad \lambda_s = \frac{\mu - \sqrt{\mu^2 + 4}}{2} < 0 \quad (\text{Rama estable})
\]

\textbf{Análisis en el equilibrio interior $P_2(1, 0)$:}
\[
J(1, 0) = \begin{pmatrix} 0 & 1 \\ -1 & \mu + 1 \end{pmatrix}
\]
\[
\det(J) = 1 > 0, \qquad \mathrm{Tr}(J) = \mu + 1
\]
\[
\begin{cases}
\mu < -1 \implies \mathrm{Tr}(J) < 0 \implies \boxed{P_2(1, 0) \text{ es un Foco/Nodo Asintóticamente Estable}} \\
\mu = -1 \implies \mathrm{Tr}(J) = 0 \implies \boxed{\text{Bifurcación de Hopf Local}} \\
\mu > -1 \implies \mathrm{Tr}(J) > 0 \implies \boxed{P_2(1, 0) \text{ es un Foco/Nodo Inestable}}
\end{cases}
\]

\vspace{3.5mm}

\textbf{\large 2. Sistema conservativo hamiltoniano no perturbado ($\mu = 0$ y sin término $xy$)}
\[
\dot{x} = y, \qquad \dot{y} = x - x^2
\]
\[
H(x, y) = \frac{1}{2}y^2 - \frac{1}{2}x^2 + \frac{1}{3}x^3 = E = \text{cte}
\]

\textbf{Ecuación analítica del lazo homoclínico no perturbado ($H = 0$):}
\[
\frac{1}{2}y^2 - \frac{1}{2}x^2 + \frac{1}{3}x^3 = 0 \iff y^2 = x^2\left(1 - \frac{2}{3}x\right)
\]
\[
\boxed{y_0(x) = \pm x\sqrt{1 - \frac{2}{3}x}, \qquad x \in \left[0, \, \frac{3}{2}\right]}
\]
\[
x_{\text{vértice}} = \frac{3}{2} = 1.5, \qquad y_{\text{máx}} = \pm\frac{1}{\sqrt{3}} \approx \pm 0.577 \quad (\text{en } x = 1)
\]

\vspace{3.5mm}

\textbf{\large 3. Cálculo de la integral de perturbación de Melnikov}
\[
g(x, y; \mu) = \mu y + xy = y(\mu + x)
\]
\[
M(\mu) = \int_{-\infty}^{\infty} y_0(t) \cdot g(x_0(t), y_0(t); \mu)\, dt = \int_{-\infty}^{\infty} y_0^2(t)\left[\mu + x_0(t)\right] dt
\]
\[
M(\mu) = \mu \int_{-\infty}^{\infty} y_0^2(t)\, dt + \int_{-\infty}^{\infty} x_0(t) y_0^2(t)\, dt = \mu I_1 + I_2
\]
\[
I_1 = \oint_{\Gamma_0} y\, dx = 2\int_0^{3/2} x\sqrt{1 - \frac{2}{3}x}\, dx = \frac{6}{5}
\]
\[
I_2 = \oint_{\Gamma_0} xy\, dx = 2\int_0^{3/2} x^2\sqrt{1 - \frac{2}{3}x}\, dx = \frac{6}{7}
\]
\[
M(\mu) = \frac{6}{5}\mu + \frac{6}{7} = 0 \implies \mu_c = -\frac{6/7}{6/5} = -\frac{5}{7}
\]
\[
\boxed{\mu_c = -\frac{5}{7} \approx -0.71428}
\]

\vspace{3.5mm}

\textbf{\large 4. Mecanismo de bifurcación homoclínica y divergencia del periodo}
\[
\begin{cases}
\mu < \mu_c \implies M(\mu) < 0 \implies W^u(P_1) \text{ no reconecta con } W^s(P_1) \text{ (flujo espirala hacia adentro)} \\
\mu = \mu_c \implies M(\mu_c) = 0 \implies \boxed{W^u(P_1) \equiv W^s(P_1) \quad (\text{Conexión Homoclínica Global})} \\
\mu > \mu_c \implies M(\mu) > 0 \implies \boxed{\text{Nace un Ciclo Límite Periódico Estable de Gran Amplitud}}
\end{cases}
\]

\textbf{Divergencia logarítmica universal del periodo al aproximarse a la bifurcación:}
\[
\lambda_u(\mu_c) = \frac{\mu_c + \sqrt{\mu_c^2 + 4}}{2} \approx \frac{-0.714 + \sqrt{0.510 + 4}}{2} \approx 0.706 > 0
\]
\[
\boxed{T(\mu) \approx -\frac{1}{\lambda_u(\mu_c)} \ln|\mu - \mu_c| + C \implies \lim_{\mu \to \mu_c^+} T(\mu) = +\infty}
\]

\vspace{3.5mm}

\textbf{\large 5. Retrato de fase de la conexión homoclínica y ciclo límite}

\begin{center}
\begin{tikzpicture}[xscale=2.2, yscale=2.0]
    % Ejes coordenados
    \draw[->, thick, gray!70] (-0.6, 0) -- (2.0, 0) node[right, black, font=\footnotesize] {$x$};
    \draw[->, thick, gray!70] (0, -1.0) -- (0, 1.05) node[above, black, font=\footnotesize] {$y = \dot{x}$};

    % Graduaciones
    \draw (1, 0.04) -- (1, -0.04) node[below, font=\scriptsize] {$1$};
    \draw (1.5, 0.04) -- (1.5, -0.04) node[below right, font=\tiny] {$1.5$};
    \draw (0.03, 0.577) -- (-0.03, 0.577) node[left, font=\tiny] {$\frac{1}{\sqrt{3}}$};
    \draw (0.03, -0.577) -- (-0.03, -0.577) node[left, font=\tiny] {$-\frac{1}{\sqrt{3}}$};

    % Lazo homoclínico crítico para mu = mu_c (rojo oscuro grueso)
    \draw[line width=1.6pt, red!80!black, domain=0:1.5, samples=16, smooth] 
        plot (\x, {\x * sqrt(max(0, 1 - (2.0/3.0)*\x))});
    \draw[line width=1.6pt, red!80!black, domain=0:1.5, samples=16, smooth] 
        plot (\x, {-\x * sqrt(max(0, 1 - (2.0/3.0)*\x))});

    % Flechas sobre el lazo homoclínico
    \draw[-{Latex[length=2.5mm]}, red!80!black, very thick] (0.85, 0.56) -- (1.15, 0.55);
    \draw[-{Latex[length=2.5mm]}, red!80!black, very thick] (1.15, -0.55) -- (0.85, -0.56);

    % Ciclo límite periódico nacido para mu > mu_c (línea continua verde)
    \draw[line width=1.3pt, solborder] 
        (0.08, 0.06) to[out=40, in=170] (0.9, 0.68) to[out=-10, in=90] (1.65, 0)
        to[out=-90, in=10] (0.9, -0.68) to[out=190, in=-40] (0.08, -0.06) to[out=140, in=-140] (0.08, 0.06);
    \draw[-{Latex[length=2.2mm]}, solborder, thick] (1.0, 0.68) -- (1.2, 0.66);
    \draw[-{Latex[length=2.2mm]}, solborder, thick] (1.2, -0.66) -- (1.0, -0.68);

    % Variedad inestable saliendo hacia la izquierda para mu < mu_c
    \draw[thick, dashed, darkpurple, -{Latex[length=2.0mm]}] 
        (0, 0) to[out=-150, in=45] (-0.45, -0.45);

    % Puntos de equilibrio
    % Punto Silla en (0, 0)
    \draw[black, thick, fill=white] (0, 0) circle (3pt);
    \draw[red!80!black, thick] (-0.07, -0.07) -- (0.07, 0.07);
    \draw[red!80!black, thick] (-0.07, 0.07) -- (0.07, -0.07);
    \node[above left, font=\scriptsize\bfseries, red!80!black] at (0, 0.08) {$P_1(0, 0)$ Silla};

    % Foco en (1, 0)
    \filldraw[unmsmblue] (1, 0) circle (2.4pt);
    \node[above, font=\scriptsize\bfseries, unmsmblue] at (1, 0.08) {$P_2(1, 0)$};

    % Etiquetas descriptivas
    \node[above right, font=\footnotesize\bfseries, red!80!black] at (1.1, 0.42) {Lazo Homoclínico ($\mu = \mu_c$)};
    \node[right, font=\footnotesize\bfseries, solborder] at (1.66, 0.15) {Ciclo Límite ($\mu > \mu_c$)};
    \node[draw=solborder, fill=white, rounded corners=0.8mm, font=\tiny\bfseries, align=left] at (1.2, -0.85) {
        $\bullet$ Conexión homoclínica en $\mu_c \approx -5/7$\\
        $\bullet$ Periodo: $T(\mu) \sim -\ln|\mu - \mu_c| \to \infty$
    };
\end{tikzpicture}
\end{center}

\end{solucion}

\begin{conclusion}
La bifurcación homoclínica es un fenómeno global no local:
\begin{enumerate}
    \item A diferencia de la bifurcación de Hopf (donde el ciclo límite nace con radio cero y frecuencia no nula), en la bifurcación homoclínica el ciclo nace con amplitud no nula finita ($R \approx 1.5$) y con frecuencia que tiende a cero ($\omega \to 0, T \to \infty$).
    \item La divergencia logarítmica $T(\mu) \sim -\ln|\mu - \mu_c|$ refleja el frenado asintótico de las trayectorias al aproximarse al punto silla hiperbólico antes de completar cada revolución.
\end{enumerate}
\end{conclusion}

\vspace{5mm}
% =========================================================================
% TEMA 4: SISTEMAS DINÁMICOS CAÓTICOS
% =========================================================================
\section*{\color{unmsmblue} Tema 4: Sistemas Dinámicos Caóticos}

\begin{problema}[Problema 4.1]
Considere el mapa tienda (Tent Map) $f: [0, 1] \to [0, 1]$:
\[
f(x) = \begin{cases} 2x, & 0 \le x \le 1/2 \\ 2(1-x), & 1/2 < x \le 1 \end{cases}
\]
Calcule analíticamente su exponente de Lyapunov $\lambda$ y determine el tiempo característico en que dos condiciones iniciales separadas por $\delta_0 = 10^{-10}$ divergen hasta una distancia de orden 1.
\end{problema}


\begin{solucion}

\textbf{\large 1. Diferenciabilidad y cálculo de la derivada casi en todas partes}
El mapa tienda $f: [0, 1] \to [0, 1]$ se define a trozos:
\[
f(x) = \begin{cases} 2x, & 0 \le x \le 1/2 \\ 2(1-x), & 1/2 < x \le 1 \end{cases}
\]
Calculamos la derivada ordinaria en los subintervalos abiertos:
\[
f'(x) = \begin{cases} +2, & x \in \left(0, \, \frac{1}{2}\right) \\ -2, & x \in \left(\frac{1}{2}, \, 1\right) \end{cases}
\]
En la cúspide $x = 1/2$, las derivadas laterales divergen en signo ($f'_-(1/2) = +2 \neq -2 = f'_+(1/2)$). Para casi todo punto $x \in [0, 1] \setminus \{1/2\}$, el módulo de la pendiente es estrictamente uniforme:
\[
\boxed{|f'(x)| = 2 \quad \left(\forall x \neq \frac{1}{2}\right)}
\]

\vspace{3.5mm}

\textbf{\large 2. Determinación de los puntos fijos y su estabilidad}
Condición de punto fijo $f(x^*) = x^*$:
\[
\begin{cases}
x \in \left[0, \, \frac{1}{2}\right]: & 2x = x \implies \boxed{x_1^* = 0} \\
x \in \left(\frac{1}{2}, \, 1\right]: & 2(1-x) = x \implies 2 - 2x = x \implies \boxed{x_2^* = \frac{2}{3}}
\end{cases}
\]
Evaluando el multiplicador de estabilidad lineal $\mu_i = f'(x_i^*)$:
\[
\mu_1 = f'(0) = 2 \implies |\mu_1| = 2 > 1, \qquad \mu_2 = f'\left(\frac{2}{3}\right) = -2 \implies |\mu_2| = 2 > 1
\]
\[
\boxed{\text{Ambos puntos fijos } x_1^* = 0 \text{ y } x_2^* = \frac{2}{3} \text{ son repelores estrictamente inestables}}
\]

\vspace{3.5mm}

\textbf{\large 3. Cálculo analítico exacto del exponente de Lyapunov}
Por la regla de la cadena aplicada a la derivada de la $n$-ésima iteración $f^n(x_0)$:
\[
(f^n)'(x_0) = \prod_{i=0}^{n-1} f'(x_i) \implies \left|(f^n)'(x_0)\right| = \prod_{i=0}^{n-1} |f'(x_i)|
\]
El exponente de Lyapunov $\lambda(x_0)$ mide la tasa asintótica media de divergencia exponencial de trayectorias cercanas:
\[
\lambda(x_0) = \lim_{n \to \infty} \frac{1}{n} \ln \left| (f^n)'(x_0) \right| = \lim_{n \to \infty} \frac{1}{n} \sum_{i=0}^{n-1} \ln |f'(x_i)|
\]
Dado que $|f'(x_i)| = 2$ para toda órbita que no incida en el conjunto numerable de preimágenes de $1/2$:
\[
\lambda = \lim_{n \to \infty} \frac{1}{n} \sum_{i=0}^{n-1} \ln 2 = \lim_{n \to \infty} \frac{n \ln 2}{n} = \ln 2
\]
\[
\boxed{\lambda = \ln 2 \approx 0.69315 > 0}
\]
Al ser $\lambda > 0$ estrictamente positivo y uniforme en todo el dominio, el sistema exhibe \textbf{caos determinista con divergencia exponencial y sensibilidad extrema a las condiciones iniciales}.

\vspace{3.5mm}

\textbf{\large 4. Divergencia exponencial y horizonte de predictibilidad (Tiempo de Lyapunov)}
Sean dos condiciones iniciales con separación infinitesimal inicial $\delta_0 = |x_0 - y_0| = 10^{-10}$.
La separación en la iteración $n$ crece según:
\[
\delta_n \approx \delta_0 \, e^{n\lambda} = \delta_0 \, (e^{\ln 2})^n = \delta_0 \cdot 2^n
\]
Determinamos el número de iteraciones $n^*$ requeridas para que la incertidumbre microscópica escale al orden macroscópico del intervalo ($\delta_{n^*} \sim 1$):
\[
\delta_0 \cdot 2^{n^*} = 1 \iff 2^{n^*} = \frac{1}{\delta_0} = 10^{10}
\]
Aplicando logaritmo natural a ambos miembros:
\[
n^* \ln 2 = 10 \ln 10 \implies n^* = \frac{10 \ln 10}{\ln 2} \approx \frac{10(2.302585)}{0.693147} \approx 33.219
\]
\[
\boxed{n^* = \lceil 33.22 \rceil = 34 \text{ iteraciones}}
\]
Cada iteración del mapa destruye exactamente $1\text{ bit}$ de información ($\log_2 |f'| = \log_2 2 = 1$). Por consiguiente, tras $34$ iteraciones, la precisión inicial de $10$ cifras decimales ($\approx 33.2\text{ bits}$) queda completamente disipada.

\vspace{3.5mm}

\textbf{\large 5. Diagrama de telaraña (Cobweb) y dinámica caótica aperiódica}

\begin{center}
\begin{tikzpicture}[scale=4.2]
    % Ejes coordenados
    \draw[->, thick, gray!70] (-0.05, 0) -- (1.15, 0) node[right, black, font=\footnotesize] {$x_n$};
    \draw[->, thick, gray!70] (0, -0.05) -- (0, 1.15) node[above, black, font=\footnotesize] {$x_{n+1}$};

    % Caja unitaria [0,1] x [0,1]
    \draw[dashed, gray!40] (1, 0) -- (1, 1) -- (0, 1);
    \draw[dashed, gray!40] (0.5, 0) -- (0.5, 1);

    % Graduaciones
    \draw (0.5, 0.02) -- (0.5, -0.02) node[below, font=\scriptsize] {$\frac{1}{2}$};
    \draw (0.667, 0.02) -- (0.667, -0.02) node[below, font=\scriptsize] {$\frac{2}{3}$};
    \draw (1, 0.02) -- (1, -0.02) node[below, font=\scriptsize] {$1$};
    \draw (0.02, 1) -- (-0.02, 1) node[left, font=\scriptsize] {$1$};
    \draw (0.02, 0.667) -- (-0.02, 0.667) node[left, font=\scriptsize] {$\frac{2}{3}$};

    % Diagonal identidad y = x
    \draw[thick, dashed, gray!80] (0, 0) -- (1.05, 1.05) node[right, font=\scriptsize, gray!90] {$y = x$};

    % Mapa tienda f(x)
    \draw[line width=1.5pt, deepteal] (0, 0) -- (0.5, 1) -- (1, 0);
    \node[above, font=\footnotesize\bfseries, deepteal] at (0.35, 0.95) {$f(x)$};

    % Trayectoria de telarana (Cobweb) con x0 = 0.18
    \draw[thick, red!85!black, -{Latex[length=1.8mm]}] (0.18, 0) -- (0.18, 0.36);
    \draw[thick, red!85!black, -{Latex[length=1.8mm]}] (0.18, 0.36) -- (0.36, 0.36);
    \draw[thick, red!85!black, -{Latex[length=1.8mm]}] (0.36, 0.36) -- (0.36, 0.72);
    \draw[thick, red!85!black, -{Latex[length=1.8mm]}] (0.36, 0.72) -- (0.72, 0.72);
    \draw[thick, red!85!black, -{Latex[length=1.8mm]}] (0.72, 0.72) -- (0.72, 0.56);
    \draw[thick, red!85!black, -{Latex[length=1.8mm]}] (0.72, 0.56) -- (0.56, 0.56);
    \draw[thick, red!85!black, -{Latex[length=1.8mm]}] (0.56, 0.56) -- (0.56, 0.88);
    \draw[thick, red!85!black, -{Latex[length=1.8mm]}] (0.56, 0.88) -- (0.88, 0.88);
    \draw[thick, red!85!black, -{Latex[length=1.8mm]}] (0.88, 0.88) -- (0.88, 0.24);
    \draw[thick, red!85!black, -{Latex[length=1.8mm]}] (0.88, 0.24) -- (0.24, 0.24);
    \draw[thick, red!85!black, -{Latex[length=1.8mm]}] (0.24, 0.24) -- (0.24, 0.48);
    \draw[thick, red!85!black] (0.24, 0.48) -- (0.48, 0.48);

    % Puntos fijos inestables marcados
    \draw[fill=white, draw=unmsmblue, thick] (0, 0) circle (1.8pt);
    \node[below left, font=\scriptsize\bfseries, unmsmblue] at (0, 0) {$x_1^* = 0$};

    \draw[fill=white, draw=unmsmblue, thick] (0.667, 0.667) circle (1.8pt);
    \node[above left, font=\scriptsize\bfseries, unmsmblue] at (0.667, 0.67) {$x_2^* = \frac{2}{3}$};

    % Anotacion condicion inicial
    \filldraw[red!85!black] (0.18, 0) circle (1.2pt);
    \node[below, font=\scriptsize\bfseries, red!85!black] at (0.18, -0.02) {$x_0 = 0.18$};

    % Caja informativa
    \node[draw=solborder, fill=white, rounded corners=0.8mm, font=\tiny\bfseries, align=left] at (0.82, 0.12) {
        $\bullet$ Pendiente: $|f'(x)| = 2$\\
        $\bullet$ Exponente: $\lambda = \ln 2 > 0$\\
        $\bullet$ Caos determinista
    };
\end{tikzpicture}
\end{center}

\end{solucion}

\begin{conclusion}
El mapa tienda ilustra el paradigma elemental del caos determinista mediante el mecanismo de estiramiento y plegado:
\begin{enumerate}
    \item La tasa uniforme de expansión $|f'(x)| = 2$ produce un exponente de Lyapunov positivo constante $\lambda = \ln 2$, coincidiendo con la entropía métrica de Kolmogorov-Sinai $h_{KS} = \ln 2$.
    \item La pérdida de predictibilidad es intrínseca y absoluta: no proviene del ruido externo, sino de la amplificación exponencial $\delta_n = \delta_0 2^n$, que restringe el horizonte determinista a un tiempo finito $n^* \propto \ln(1/\delta_0)$.
\end{enumerate}
\end{conclusion}

\vspace{3mm}

\begin{problema}[Problema 4.2]
Para el mapa logístico $x_{n+1} = r x_n(1 - x_n)$ con $x \in [0, 1]$:
Determine analíticamente el valor crítico $r_1$ donde la órbita de periodo 1 pierde estabilidad y nace la órbita de periodo 2.
\end{problema}

\vspace{5mm}

\begin{problema}[Problema 4.3]
La transición al caos en mapas unimodales sigue una cascada infinita de duplicaciones de periodo $2^n$ que ocurren en parámetros $r_n$. Si $r_1 = 3$ y $r_2 = 1 + \sqrt{6} \approx 3.44949$, estime el umbral crítico de acumulación caótica $r_\infty$ usando la constante de Feigenbaum $\delta \approx 4.6692016$.
\end{problema}

\vspace{5mm}

\begin{problema}[Problema 4.4]
El sistema de Lorenz (1963) está dado por:
\[
\dot{x} = \sigma(y - x), \qquad \dot{y} = r x - y - x z, \qquad \dot{z} = x y - b z
\]
con $\sigma, r, b > 0$.
\begin{enumerate}[label=\textbf{\alph*)}]
    \item Demuestre que el flujo es estrictamente disipativo y calcule la tasa de contracción de volumen en el espacio de fases.
    \item Demuestre que todas las trayectorias entran eventualmente y permanecen atrapadas en un elipsoide acotado de $\mathbb{R}^3$.
\end{enumerate}
\end{problema}

\vspace{5mm}

\begin{problema}[Problema 4.5]
Demuestre que para cualquier flujo suave tridimensional disipativo $\dot{\mathbf{x}} = \mathbf{f}(\mathbf{x})$ que posea un atractor caótico acotado que no contenga puntos de equilibrio, el espectro de exponentes de Lyapunov $(\lambda_1, \lambda_2, \lambda_3)$ debe satisfacer rigurosamente:
\[
\lambda_1 > 0, \qquad \lambda_2 = 0, \qquad \lambda_3 < 0, \qquad \text{con } \lambda_1 + \lambda_2 + \lambda_3 < 0
\]
\end{problema}

\vspace{5mm}
% =========================================================================
% TEMA 5: ATRACTORES EXTRAÑOS Y GEOMETRÍA FRACTAL
% =========================================================================
\section*{\color{unmsmblue} Tema 5: Atractores Extraños y Geometría Fractal}

\begin{problema}[Problema 5.1]
El corte transversal de la mayoría de atractores extraños posee estructura de conjunto de Cantor. Calcule rigurosamente la dimensión de caja (Box-Counting Dimension) $D_0$ del conjunto ternario clásico de Cantor.
\end{problema}

\vspace{5mm}

\begin{problema}[Problema 5.2]
Describa el mecanismo geométrico de la herradura de Smale y explique por qué genera un atractor extraño con infinitas órbitas periódicas inestables y dinámica caótica simbólica.
\end{problema}

\vspace{5mm}

\begin{problema}[Problema 5.3]
Considere el mapa bidimensional de Hénon (1976):
\[
x_{n+1} = 1 - a x_n^2 + y_n, \qquad y_{n+1} = b x_n
\]
\begin{enumerate}[label=\textbf{\alph*)}]
    \item Calcule el determinante de su matriz jacobiana y verifique la contracción de áreas en el plano para los parámetros canónicos $a = 1.4$ y $b = 0.3$.
    \item Demuestre que el mapa es invertible y obtenga su mapeo inverso explícito $T^{-1}(x, y)$.
\end{enumerate}
\end{problema}

\vspace{5mm}

\begin{problema}[Problema 5.4]
El sistema de Rössler (1976) viene dado por:
\[
\dot{x} = -y - z, \qquad \dot{y} = x + a y, \qquad \dot{z} = b + z(x - c)
\]
con $a = b = 0.2$ y $c = 5.7$.
Explique cómo el flujo tridimensional se reduce a un mapa unidimensional con cúspide a través de la sección de Poincaré transversal $y = 0, \dot{y} > 0$.
\end{problema}

\vspace{5mm}

\begin{problema}[Problema 5.5]
La Conjetura de Kaplan-Yorke (dimensión de Lyapunov $D_L$) relaciona el espectro de exponentes de Lyapunov ordenados $\lambda_1 \ge \lambda_2 \ge \dots \ge \lambda_n$ con la dimensión de información del atractor mediante:
\[
D_L = k + \frac{\sum_{j=1}^k \lambda_j}{|\lambda_{k+1}|}
\]
donde $k$ es el entero más grande tal que $\sum_{j=1}^k \lambda_j \ge 0$.
Para el atractor de Lorenz con parámetros canónicos $\sigma = 10, b = 8/3, r = 28$, los exponentes calculados numéricamente son:
\[
\lambda_1 \approx 0.9056, \qquad \lambda_2 = 0, \qquad \lambda_3 \approx -14.5723
\]
\begin{enumerate}[label=\textbf{\alph*)}]
    \item Compruebe analíticamente que $\lambda_1 + \lambda_2 + \lambda_3$ coincide con la divergencia media del sistema de Lorenz.
    \item Calcule la dimensión fractal de Lyapunov $D_L$ del atractor de Lorenz e interprete físicamente su valor fraccionario.
\end{enumerate}
\end{problema}

\vspace{5mm}

\end{document}
